% 本文件由 LaTeX 排版 Agent 根据有效 translation.json 自动生成。
% author=John Chrysostom; work_id=2307; title=弟茂德后书讲道集

\chapter{弟茂德后书讲道集}

% structure: p0001 source=h1 target=omitted reason=page-title-already-rendered-by-parent

讲道一\newline{}讲道二\newline{}讲道三\newline{}讲道四\newline{}讲道五\newline{}讲道六\newline{}讲道七\newline{}讲道八\newline{}讲道九\newline{}讲道十\par

\SourceRecord{Translated by Philip Schaff.；Nicene and Post-Nicene Fathers, First Series；Vol. 13.；Edited by Philip Schaff.；Buffalo, NY: Christian Literature Publishing Co.,；1889.；<http://www.newadvent.org/fathers/2307.htm>.}

% structure: p0001 source=h1 target=section reason=page-title

\section{弟茂德后书讲道词 第一讲}

\BibleRef{弟后 1:1-2}\par

\BibleQuote{保禄，因天主的旨意，为着在基督耶稣内生命的恩许，作耶稣基督的宗徒，致书给蒙爱的儿子弟茂德：愿恩宠、怜悯与平安，由圣父和我们的主基督耶稣赐与你。}\BibleRef{弟后 1:1-2}\par

他写这第二封致弟茂德书的原因是什么？他曾说过：\BibleQuote{我希望能很快到你那里去}\BibleRef{弟前 3:14}，但这并未实现，他并未到弟茂德那里去，反而用一封信来安慰他，因为弟茂德或许正因他的缺席而悲伤，并因他现已承担的治理之重任而感到压力。连伟大的人物，当他们在舵位上，肩负着管理教会的职责时，也会感到其职位的陌生，并为事务的波涛所淹没，犹如被压垮一般。这尤其在福音初次传扬时更为真实，那时处处土地未经开垦，周遭尽是敌对与敌意。此外，还有从犹太教师兴起的异端，正如他在前书中所指出的。他不仅以书信安慰弟茂德，还邀请他到自己这里来：他说：\BibleQuote{你要赶紧到我这里来，}以及\BibleQuote{你来的时候，要把书带来，尤其是那些羊皮卷。}\BibleRef{弟后 4:9}看来他写这封信时，他的结局已经临近。因为他说：\BibleQuote{我现在已被奠祭}；又说到：\BibleQuote{在我初次申诉时，没有一人站在我身旁。}\BibleRef{弟后 4:6}为纠正这一切，他既以自己的试炼提供安慰，又说：\par

\BibleQuote{保禄，因天主的旨意，为着在基督耶稣内生命的恩许，作耶稣基督的宗徒。}\BibleRef{弟后 1:1}\par

这样，他在开头就提升了他的心灵。不要对我讲这里的危险，他说。这些危险为我们赢得永生，那里没有危险，只有忧伤和悲哀的消逝。因为他使我们成为宗徒，不仅是为了面对危险，更是为了受苦难甚至死亡。而如果向他叙述自己的苦难并不能安慰他，反而增加他的忧愁，那么他就立刻以安慰开始，说：\BibleQuote{按着在基督耶稣内生命的恩许。}如果它是\Quote{恩许}，就不要在此世寻求它。因为\BibleQuote{人所看见的，不是希望。}\BibleRef{罗 8:24}\par

\BibleQuote{致蒙爱的儿子弟茂德。}\BibleRef{弟后 1:2}\par

不仅是他的\Quote{儿子}，而是\Quote{蒙爱的}；因为儿子也可能不被爱。他说，你不是这样；我称你不仅为儿子，而是\Quote{蒙爱的儿子}。正如他称迦拉达人为他的儿女，但同时也责备他们：他说：\BibleQuote{我的孩子们，我为你们再受产痛。}\BibleRef{迦 4:19}他特别以称他为蒙爱的来见证他的德行。因为那里爱不是出于本性，就必然出于对象的功德。生自我们的人，被爱不仅因他们的德行，也因本性的力量；但那些因信仰而被爱的人被爱则只是因他们的功德，此外还能有什么原因呢？尤其是对保禄而言，他从不偏私行事。而且，他称他为\Quote{蒙爱的儿子}，也表明他并非因恼怒他、轻视他或谴责他才没有到他那里去。\par

\BibleQuote{愿恩宠、怜悯与平安，由圣父和我们的主基督耶稣赐与你。}\BibleRef{弟后 1:2}\par

他在开始时就祈求的这些事，现在再次为他祈求。注意，他一开始就为自己没有到他那里去、没有见他而辩解。因为他说：\Quote{直到我来}，以及\Quote{希望很快到你那里去}，这些话曾使弟茂德期待他很快到来。为此他辩解，但没有立即说明不来的原因，免得使他过度悲伤。因为他被皇帝囚禁在监狱里。但在他写信末尾邀请他到他自己那里去时，他才告诉他。他并没有一开始就使他陷入悲伤，而是鼓励他希望能见到他。\BibleQuote{迫切渴望见到你}，以及\BibleQuote{你要赶紧到我这里来。}\BibleRef{弟后 1:4}于是他立刻振作他，并称赞他。\par

第3、4节。\BibleQuote{我感谢天主，就是我以纯洁的良心所事奉的，像我祖先一样，我不断在日夜的祈祷中记念你；迫切渴望见到你，记念你的眼泪，使我满心喜乐。}\BibleRef{弟后 1:3-4}\par

\Quote{\InnerQuote{我感谢天主}，他说，\InnerQuote{我记念你}，我是如此爱你。}这是极爱的标志，当一个人因深爱而夸耀自己的情感时。\Quote{我感谢天主}，他说，\Quote{我所事奉的}：是怎样的呢？\Quote{以纯洁的良心}，因为他没有违背自己的良心。此处他谈到自己无可指摘的生活，因为他处处将生活称为良心。或者因为无论为了任何人的理由，我从未放弃任何我意图的善事，甚至在我作迫害者时也是如此。因此他说，\Quote{我蒙了怜悯，因为我是在不信中无知而做的} \BibleRef{弟前 1:13}；几乎是说，\Quote{不要猜测那是出于邪恶。}他恰当地称赞自己的品性，使他的爱显得真诚。因为他所说的其实是，\Quote{我不虚假，我的心口如一。}同样在\WorkTitle{宗徒大事录}中，我们读到他是被迫称赞自己。因为当他们诽谤他是叛乱者和革新者时，他为自己辩护说，\Quote{\Person{阿纳尼雅}对我说：我们祖先的天主拣选了你，叫你知道他的旨意，看见那位义者，并听见他口中的声音。因为你将为他作见证，向众人见证你所看见所听见的。} \BibleRef{宗 22:14}同样在此，为了不让人以为他善忘，而具有缺乏友誼和良心的品性，他公正地称赞自己，说，\Quote{我无时无刻不记念你}，而且不仅如此，\Quote{在我的祈祷中。}这就是说，这是我祈祷的事务，是我不断致力的事。因为他用这话表明，\Quote{为此我日夜恳求天主，渴望见到你。}注意他的热切渴望，他爱的强烈。还有他的谦卑，他如何向门徒道歉，然后表明这不是出于轻浮或虚浮的理由；这是他以前向我们显明的，现在再次证明。\Quote{记念你的眼泪。}\Person{弟茂德}在与他分离时伤心流泪，这很自然，比一个婴儿被从母乳和母亲怀中夺走更甚。\Quote{使我充满喜乐；极切渴望见到你。}我不愿自愿剥夺自己这么大的快乐，即使我生性无情粗暴，那些眼泪进入我的记忆也足以使我软化。但我并非如此性格。我是那些纯洁事奉天主的人之一；因此许多强烈的动机催逼我来见你。于是他哭了。他还提到另一个原因，那是安慰性的。\par

第5节。\Quote{当我记起你无伪的信德。} \BibleRef{弟后 1:5}\par

这是另一项称赞，即\Person{弟茂德}不是来自外邦人，也不是来自不信者，而是来自一个从一开始就事奉基督的家庭。\BibleRef{宗 16:1}\par

\Quote{这信德曾先住在你的祖母\Person{洛依斯}和你的母亲\Person{欧尼刻}心里。} \BibleRef{弟后 1:5}\par

关于\Person{弟茂德}，经上说，\Quote{他是某个犹太妇人的儿子，并且信了主。}怎么是犹太妇人？怎么信了？因为她不是外邦人，\Quote{但因他父亲是希腊人，又因那地方的犹太人，他就带他行了割礼。} \BibleRef{宗 16:1}\par

\Quote{我也深信这信德也在你心里。} \BibleRef{弟后 1:5}\par

他的意思是，你从一开始就有这卓越之处。你从祖先领受了无伪的信德。因为我们祖辈的赞誉，当我们有分时，也归于我们。否则它们毫无用处，反而谴责我们；因此他说，\Quote{我也深信这信德也在你心里。} \BibleRef{弟后 1:5}这不是猜测，他说，这是我的确信；我对此完全肯定。因此，如果你不是出于任何人的动机而接受了它，没有什么能动摇你的信德。\par

第6节。\Quote{为此我提醒你，激发天主的恩赐，即因我按手在你心里所赋予的。} \BibleRef{弟后 1:6}\par

你看他以为\Person{弟茂德}是何等沮丧消沉。他几乎说，\Quote{不要以为我轻视你，但要确信我不定你的罪，也没有忘记你。无论如何，想想你的母亲和祖母。正是因为我知你有无伪的信德，我才提醒你。}因为激发天主的恩赐需要很大热忱。如火需要燃料，同样恩宠需要我们的殷勤，才能常保持炽热。\Quote{我提醒你，激发天主的恩赐，即因我按手在你心里所赋予的，}那就是，你所领受的圣神的恩宠，为管理教会、行奇迹和各样服务。因为这恩宠，我们有权点燃或熄灭；因此他别处说，\Quote{不要熄灭圣神。} \BibleRef{得前 5:19}因为因懈怠和疏忽它被熄灭，因警醒和勤勉它得以保持。它确实在你内，但你要使它更加旺盛，即充满信心、喜乐和欢愉。要刚强地站立。\par

第7节。\Quote{因为天主赐给我们的，不是胆怯的精神，而是大能、爱德和健全理性的精神。} \BibleRef{弟后 1:7}\par

这就是说，我们领受圣神，不是为了退缩逃避劳苦，而是为了能放胆宣讲。因为在列王纪的战事中，祂把恐惧之神赐给许多人。\BibleQuote{恐惧之神降在他们身上。}\BibleRef{出 15:16}这就是说，祂把恐怖倾注在他们身上。相反地，祂赐给你们的，却是大能之神和爱德之神。这固然是出于恩宠，却又不纯粹是恩宠，而是当我们先履行了自己的本分之后。因为那使我们呼喊\Quote{阿爸，父啊}的圣神，激发我们对祂和近人的爱德，好使我们彼此相爱。因为爱德由能力和不惧怕而生。因为没有什么比恐惧和猜忌更容易摧毁爱德。\par

\Quote{因为天主赐给我们的，不是恐惧之神，而是大能之神、爱德之神和慎重之神}：他把灵魂的健康状态称为慎重之心态，或者可能指心灵的清醒，或者心灵的镇定，好使我们保持清醒，如果有什么灾祸降临到我们身上，它能让我们冷静，并除去多余之物。\par

道德训诲。那么，我们不要为临到我们身上的灾祸而苦恼。这就是心灵的清醒。他说：\Quote{在试探的季节，不可急躁。}\BibleRef{德 2:2}许多人在家中各有各的忧伤，我们分担彼此的痛苦，尽管并非分担其根源。因为有人因妻子而不幸，有人因孩子或仆人，有人因朋友，有人因仇敌，有人因邻人，有人因某种损失。忧伤的原因多种多样，所以我们找不到任何人完全没有某种烦恼和不幸，只是有些人忧伤更大，有些人较小。因此，我们不要不耐烦，也不要以为只有我们自己是不幸的。\par

因为在这必死的生命中，没有什么是免于忧愁的。若非今日，便是明日；若非明日，便是某日，烦恼必然来临。正如一个人不能航行，我的意思是，在广阔的大海上而丝毫不觉不安，同样，他也无法度过此生而不经历忧愁，即使你说的是富人；因为他富有，就有许多放纵欲望的机会，即使是国王本人，因为他也要受许多人管辖，不能为所欲为。他被迫违背自己的意愿赐予许多恩惠，而且比所有人都更不得不做他不愿意做的事。何以如此？因为他身边有许多人希望得到他的礼物。想想当他渴望做成某事，却因恐惧或猜疑，或被敌人或朋友阻碍而无法实现时，他会有多么沮丧。常常当他成功达成某个目标时，却因许多人成为他的敌人而失去所有乐趣。再者，你认为那些过着安逸生活的人就没有忧愁吗？这不可能。正如人不能逃避死亡，他也不能逃避忧愁。他们必须忍受多少烦恼，我们无法用言语表达，只有他们自己才能亲身体验！多少人在他们的财富和奢侈中千百次祈求死亡！因为奢侈绝不能使人远离忧愁：它恰恰是产生忧伤、疾病和不安的原因，常常在根本没有真实理由的情况下。因为当灵魂处于此种状态时，它甚至会在没有原因的情况下悲伤。医生说，胃的虚弱状态会无缘无故地产生忧伤；难道我们自己不也发生同样的事吗？感到不安，却不知为何？简而言之，我们找不到任何一个人是免于忧愁的。即使他的忧愁少于我们，他也并不这样认为，因为他感受到自己的忧愁，胜过别人的。正如身体某处疼痛的人，认为自己的痛苦超过邻人。患眼疾的人，认为没有什么比这更痛苦；患胃病的人，认为这是最严重的疾病；每个人都认为自己所受的痛苦最沉重。忧愁也是如此，每个人都认为当前的忧愁最为严重。因为他根据自己的经历来判断。没有孩子的人认为没有比没有孩子更悲伤的事；贫穷且子女众多的人，抱怨大家庭的极端痛苦；只有一个孩子的人，把这看作最大的不幸，因为那唯一的孩子，由于过分珍视，从不纠正，变得任性，给父亲带来忧伤；妻子漂亮的人，认为没有什么比拥有漂亮的妻子更糟糕，因为这是嫉妒和阴谋的缘由；妻子丑陋的人，认为没有什么比相貌平平的妻子更糟糕，因为这始终令人不快。平民认为没有什么比他的生活方式更卑微、更无用的了。士兵宣称没有比战争更劳累、更危险的了；他宁愿靠面包和水生活，也不愿忍受那样的艰辛。掌权者认为没有比照顾他人的需要更沉重的负担。受权者认为没有比仰人鼻息更奴性的事了。已婚者认为没有什么比妻子和婚姻的烦恼更糟的了。未婚者宣称没有比不结婚、没有家庭的安宁更可怜的了。商人认为农夫因安全而幸福。农夫认为商人因财富而幸福。简而言之，全人类都不知何故难以满足，不满且急躁。当谴责整个人类时，他说：\Quote{人是虚无之物} \BibleRef{咏 143:4}，暗示整个人类是可怜不幸的受造物。有多少人渴望老年！有多少人认为青年是幸福的时光！因此每个不同的时期都有其不幸。当我们因年轻而受责备时，我们会说：为什么我们不是老年人？而当我们的头发灰白时，我们又问：我们的青春飞到哪里去了？总之，忧愁的机会数不胜数。只有一条道路可以逃避这种起伏。那就是美德的道路。然而它也有自己的忧愁，只是这些忧愁并非无益，而是带来收获和益处。因为如果有人犯了罪，他因忧愁而来的痛悔洗去他的罪。或者，如果他因同情堕落的弟兄而悲伤，这并非没有报酬。因为同情患难中的人，使我们在天主面前获得极大的信心。\par

因此请听，从圣经中\Person{约伯}的例子所学到的智慧！也请听\Person{保禄}所说的：\Quote{与哭泣的人一同哭泣}；又说：\Quote{俯就卑微的人。}\BibleRef{罗 12:15}因为，通过分担忧愁，它的重担就减轻了。正如在沉重的负荷下，分担一部分重量的人减轻了独自承受者的负担，在其他一切事上也是如此。\par

但如今，当我们的一位亲人去世时，有许多人坐在旁边安慰我们。甚至一头驴跌倒了，我们也常常把它扶起来；但当我们的弟兄的灵魂跌倒时，我们却忽略它们，从旁走过，仿佛它们不如一头驴有价值。如果我们看见有人不体面地走进酒馆，甚至看见他喝醉，或犯任何其他不雅行为，我们不制止他，反而与他同流。因此\Person{保禄}说：\BibleQuote{他们不仅做这些事，还赞同别人做。}\BibleRef{罗 1:32}大多数人甚至为醉酒的目的而结社。但是，人啊，你要结社以制止酗酒的疯狂。这样的友爱行动有益于那些被捆绑或受苦的人。\Person{保禄}曾吩咐格林多人这样的事，暗示此事时说：\BibleQuote{免得我来的时候再有聚集。}\BibleRef{格前 16:2}但现在一切为了奢侈、狂欢和享乐。我们有共同的座位、共同的桌子、共同的酒、共同的开销，却没有共同的施舍。宗徒时代就有这样的友爱行动；他们把一切财物归入公共储备。现在我不要求你们全部奉献，只奉献一部分。\Quote{每七日的第一日，每人要照自己的进项抽出来留着，正如天主使他蒙福一样。}并把它作为七日的献金存起来。这样施舍，不论多少。\BibleQuote{因为你们不可空手朝见主。}\BibleRef{出 23:15}这话是对犹太人说的，何况对我们。为此，穷人站在门前，为叫没有人空手进入，但每人进门时都要施舍。你们进来恳求怜悯，先要显示怜悯。后进来的人欠得更多。因为我们若先施舍，后来的人就偿还更多。使天主成为你的债务人，然后献上祈祷。借给祂，然后求回报，你必连本带利收回。天主愿意这样，祂不反悔。若你以施舍祈求，祂就自觉有责。若你以施舍祈求，你就是借出并收取利息。是的，我恳求你们！你们被垂听不是因为伸开双手！伸出你们的手，不是向天，而是向穷人。如果你们伸出的手触摸到穷人的手，你们就达到了天堂的最高处。因为坐在那里的那一位接受你们的施舍。但如果你们举起空手，就一无所获。如果一位身披紫袍的国王来到你们面前求施舍，你们岂不乐意把一切给他？但现在，当你们透过穷人被恳求，不是被地上的王，而是被天上的王恳求时，你们竟然漠不关心，拖延你们的献金？那么你们该受何等惩罚呢？因为被垂听并不在于举起你们的双手，也不在于多言多语，而在于你们的善工。因为请听先知的话：\BibleQuote{你们伸出手时，我必掩目不看；你们行大祈祷时，我也不俯听。}\BibleRef{依 1:15}因为需要怜悯的人应当沉默，甚至不得仰望上天；那有自信的人才能多言。但圣经说什么呢？\BibleQuote{为孤儿伸冤，为寡妇辩护，学习行善。}\BibleRef{依 1:17}如此，我们必被垂听，即使我们不举手，不言语，也不请求。因此，让我们在这些事上热心，好获得所应许的福分，藉着恩宠和慈爱，等等。\par

\SourceRecord{Translated by Philip Schaff.；Nicene and Post-Nicene Fathers, First Series；Vol. 13.；Edited by Philip Schaff.；Buffalo, NY: Christian Literature Publishing Co.,；1889.；<http://www.newadvent.org/fathers/230701.htm>.}

% structure: p0001 source=h1 target=section reason=page-title

\section{\WorkTitle{弟茂德后书第二篇讲道}}

\BibleRef{弟后 1:8-10}\par

\BibleQuote{所以，你不要以给我们的主作证为耻，也不要以我这为主被囚的为耻；但要依赖天主的大能，为福音同我一起受苦。天主拯救了我们，以圣召召叫了我们，并不是按照我们的行为，而是按照祂的旨意和恩宠；这恩宠是在万世以前，在基督耶稣内赐给我们的，如今藉着我们的救主耶稣基督的显现，显示出来了。}\par

没有比人用人的理性来量度和判断天主的事更糟的了。因为这样做，他将从那块磐石上坠落遥远的距离，并被剥夺光明。因为如果有人想用肉眼去捕捉太阳的光芒，他不仅不能捕捉到，而且除了失败之外，还会遭受极大的伤害；同样，但程度更高，那些想用人的理性凝视那光的人，将会遭受此害，并滥用天主的恩赐。因此，请看马尔基翁、摩尼、瓦伦提诺以及其他将异端和有害教义引入天主教会的人，他们用人的理性来量度天主的事，就对天主的救世计划感到羞耻。然而，这不是可羞耻的事，而是可夸耀的；我是指基督的十字架。因为天主对人的爱，没有比这更大的标记了，不是天，不是海，不是地，不是从无中创造万物，也不是其他任何事物，而是十字架。因此，保禄夸耀说：\BibleQuote{我绝不以别的夸耀，只以我们的主耶稣基督的十字架夸耀。}\BibleRef{迦 6:14} 但属血气的人，以及那些把天主看得不比人更高的人，就绊倒并感到羞耻。因此，保禄从一开始就劝诫他的门徒，并通过他劝诫所有其他人，说：\BibleQuote{你不要以给我们的主作证为耻，} 也就是说，\Quote{不要羞于宣讲那位被钉十字架者，反而要以此为荣。} 因为就其本身而言，死亡、监禁和锁链是可耻和可辱的事。但当原因摆在我们面前，并正确看待这一奥迹时，它们将显得充满尊严，是可夸耀的。因为正是那死亡拯救了将要灭亡的世界。那死亡连接了地与天，那死亡摧毁了魔鬼的权势，使人成为天使和天主的子女：那死亡将我们的本性提升到了君王的宝座。那些锁链使许多人悔改。\BibleQuote{不可} 因此 \BibleQuote{羞愧}，他说，\BibleQuote{以我们的主作证为耻，也不要以我这被囚的为耻；但要为福音同我一起受苦}；就是说，即使你遭遇同样的事，也不要羞愧。因为这一点从上面所引的经文可以看出来：\BibleQuote{天主赐给了我们刚毅、爱德和谨慎的精神}；以及随后所说的：\BibleQuote{你要为福音同我一起受苦}：不仅是不要羞愧于这些事，而且连经历它们也不要羞愧。\par

他没有说：\Quote{不要害怕，} 而是为了更鼓励他，说：\Quote{不要羞愧，} 好像如果他克服了羞愧，就没有更大的危险了。因为只有当人被羞愧胜过时，它才令人压抑。所以，如果我现在是一个囚犯，我——曾复活死者、行奇迹、走遍世界的人——你也不要因此羞愧。因为我被囚禁，不是作为罪犯，而是为了那位被钉十字架者的缘故。如果我的主不以十字架为耻，我也不以锁链为耻。非常恰当的是，当他劝告他不要羞愧时，他提醒他十字架。他的意思是：如果你不以十字架为耻，也不要因锁链而羞愧；如果我们的主和师傅忍受了十字架，我们岂不更该忍受锁链？因为对那些祂所忍受的事感到羞耻，就是对被钉十字架者的羞耻。我现在不是为自己而戴这些锁链；所以不要顺从人的情感，而要分担这些苦难。\BibleQuote{你要为福音同我一起受苦。} 他说这话，不是好像福音能受到损害，而是为了激励他的门徒为福音受苦。\par

\BibleQuote{依照天主的大能；祂拯救了我们，以圣召召叫了我们，不是按照我们的行为，而是按照祂自己的旨意和恩宠；这恩宠是在万世以前，在基督耶稣内赐给我们的。}\par

尤其因为说\Quote{你们要分担苦难}是一件难事，他又安慰他。你要认为你忍受这些事，不是凭自己的能力，而是凭天主的大能。因为选择并热心是你的部分，而减轻苦难并使其停止是天主的事。然后祂向他显示祂大能的证明。想想你是如何得救的，你是如何被召叫的。正如他在别处所说：\BibleQuote{按照那在我们内运行的大能。}\BibleRef{弗 3:20} 说服世界相信，比创造诸天更显示了更大的大能。但他是如何\Quote{以圣召被召叫}的呢？这就是说，祂使他们成圣，他们本是罪人和仇敌。而这并不是出于我们自己，而是天主的恩赐。如果祂在召叫我们时既是全能的，又是良善的，因为祂出于恩宠而非义务做了这事，我们就不应该害怕。因为那位当我们本应灭亡时，却以恩宠拯救了我们这些仇敌的祂，难道当祂看到我们工作时，不更会与我们合作吗？他说：\BibleQuote{不是按照我们自己的行为，}他说，\BibleQuote{而是按照祂自己的旨意和恩宠，} 就是说，没有人强迫祂，没有人劝告祂，而是出于祂自己的旨意，因祂自己良善的推动，祂拯救了我们；这就是\BibleQuote{按照祂自己的旨意}的意思。\BibleQuote{这恩宠是在万世以前赐给我们的。} 也就是说，从无始以来就决定了这些事要在基督耶稣内完成。这不是一个轻率的考虑：祂从一开始就愿意。这不是事后的主意。那么，圣子怎么不是永恒的呢？因为祂也从起初就愿意了。\par

第10节。\BibleQuote{如今藉着我们的救主耶稣基督的显现，显示出来了；祂已废除了死亡，并藉着福音将生命和不朽彰显出来。}\BibleRef{弟后 1:10}\par

你看到了德能，你看到了恩赐，不是因行为赐下的，而是藉着福音。这些是希望的对象：因为二者都在祂的身体内成就。它们如何在我们内成就呢？藉着福音。\par

第11节。\BibleQuote{我为此被立为宣讲者、宗徒，以及外邦人的教师。} \BibleRef{弟后 1:11}\par

他为何如此频繁地重复这一点，并称自己为外邦人的教师？因为他愿意说服他们，他们也应亲近外邦人。所以不要因我的苦难而丧胆。死亡的筋腱已被松弛。我受苦，不是作为作恶者，而是因为我是\Quote{外邦人的教师}。同时，他使他的言论值得信赖。\par

第12节。\BibleQuote{因此，我也受这些苦，然而我不以为耻。因为我知道我所相信的是谁，也深信祂能保全我所交付给祂的，直到那一日。} \BibleRef{弟后 1:12}\par

\Quote{我不以为耻，}他说。因为锁链、苦难，难道是耻辱的事吗？那么，你们不要羞愧！你看他如何以行为阐明他的教训。\Quote{这些事，}他说，\Quote{我受苦}：我被投入监狱，我被流放；\BibleQuote{因为我知道我所相信的是谁，也深信祂能保全我所交付给祂的，直到那一日。}什么是\Quote{所交付的}？就是信德、福音的宣讲。他说，将那交付给祂的，会保全它不受损伤。我忍受一切，以免被剥夺这宝藏，只要它保存完好，我就不以这些事为耻。或者他将信友称为天主托付给他的，或他托付给天主的。因为他说：\BibleQuote{现在我把你们交托给主。} \BibleRef{宗 20:32} 也就是说，这些事对我不会是无益的。而在\Person{弟茂德}身上，可以看出这\Quote{所交付的}托付的果实。你看，由于他对门徒所怀的希望，他对苦难无动于衷。\par

训诲。教师应如此，应如此看待门徒，视他们为一切。他说：\BibleQuote{现在我们活着，}他说，\BibleQuote{如果你们在主内站稳。}他又说：\BibleQuote{我们的希望、喜乐或所夸耀的冠冕是什么呢？不是你们在我主耶稣基督面前吗？}\BibleRef{得前 3:8}你们看到他对此事的焦虑，他对门徒得救的关切不亚于自己。因为教师应当超越生身父母，比他们更热切。而门徒也当以善意待他们。因为他说：\BibleQuote{你们要服从领导你们的人，且要听命，因为他们为你们的灵魂时刻警醒，好像要交账的人。}\BibleRef{希 13:17}试问，他肩负如此危险的职责，而你们却不愿服从他，更何况这是为了你们自己的益处？因为即使他本人的状况良好，但只要你们处于恶劣状况，他的焦虑就持续，他要交双重的账。试想，为每一个属下负责和焦虑意味着什么？什么样的尊荣能与这样的危险相称，什么样的服事能回报这样的危险？你无法给出等价。因为你还没有为他献出你的灵魂，但他却为你舍命；即使在此世没有舍命，但在必要时他会在那里失去它。然而你们连在言语上也不愿服从。这就是一切祸患的首要原因：统治者的权威被忽视，没有敬意，没有恐惧。他说：\BibleQuote{你们要服从领导你们的人，且要听命。}但如今一切都颠倒混乱。我说这话不是为了统治者的缘故（因为从我们给予他们的尊荣中，他们能得到什么益处呢？除了我们变得更加顺服之外），而是为了你们的益处。因为对于将来，他们不会因所受的尊荣而得利，反而会受到更大的定罪；他们也不会因所受的恶待而受损，反而会得到更多的原谅。但这一切我都是为你们而说。因为当统治者受人民尊敬时，这也会被记在他们账上；就如厄里的事上所说：\BibleQuote{我不是从他父家拣选了他吗？}\BibleRef{撒上 2:27}但当他们受侮辱时，就如撒慕尔的事上，天主说：\BibleQuote{他们不是弃绝了你，而是弃绝了我。}\BibleRef{撒上 8:7}因此，侮辱是他们的益处，尊荣是他们的负担。所以我说这话是为你们，不是为他们。尊敬司铎的人，也必尊敬天主；学会轻视司铎的人，久而久之也会侮辱天主。祂说：\BibleQuote{谁接待你们，就是接待我。}\BibleRef{玛 10:40}又说：\BibleQuote{要尊敬司铎。}\BibleRef{德 7:31}犹太人因为轻视梅瑟而学会轻视天主，甚至要用石头砸死他。因为一个人若对司铎有虔诚之心，他对天主就更加虔诚。即使司铎是邪恶的，天主见你因敬畏祂而尊敬那虽不配受尊敬的人，祂自己必会报答你。因为如果\BibleQuote{谁因先知的名接待先知，必得先知的赏报}\BibleRef{玛 10:41}，那么谁尊敬、服从、听从司铎，必得赏报。因为若在款待上，当你不知客人是谁时，竟得到如此高的报酬，何况你服从祂命你服从的人呢？祂说：\BibleQuote{经师和法利塞人坐在梅瑟的座位上；凡他们所吩咐你们的，你们都要谨守遵行；但不要效法他们的行为。}\BibleRef{玛 23:2}难道你们不知道司铎是什么吗？他是上主的天使。他说的是自己的话吗？如果你们轻视他，你们轻视的不是他，而是立他为司铎的天主。但你们可能会问：怎能看出他是天主所立的？如果你们以为不然，那么你们的希望就落空了。因为如果天主不藉着他行事，你们就没有洗礼，不领受奥迹，不领受祝福的恩惠；那么你们就不是基督徒。你们会说：那么天主立了所有人，甚至不配的人？天主固然不立所有人，但祂藉着所有人行事，即使他们本身不配，为使人民得救。因为如果祂为人民的缘故，曾藉着一头驴和非常邪恶的巴郎说话，那么祂更会藉着司铎的口说话。为了我们的得救，天主有什么不做、不说的呢？祂不藉着谁行事呢？因为如果祂曾藉着犹达斯和那些\BibleQuote{说预言}的人行事（祂将对他们说：\BibleQuote{我从来不认识你们；你们这些作恶的人，离开我吧！}\BibleRef{玛 7:22}），如果其他人\BibleQuote{驱逐魔鬼}\BibleRef{咏 6:8}，那么祂岂不更藉着司铎行事？因为如果我们查考统治者的生活，我们就成了我们自己教师的任命者，一切都将混乱：脚在上，头在下。请听保禄说：\BibleQuote{对我来说，受你们审断，或受人间法庭审断，都是极小的事。}\BibleRef{格前 4:3}又说：\BibleQuote{你为什么判断你的弟兄？}\BibleRef{罗 14:10}因为如果我们不能判断我们的弟兄，何况我们的教师呢？如果天主命令这样行，你们行的好，不行便是犯罪；但如果相反，你们切勿这样做，也不要试图越界。亚郎造了金牛犊之后，科辣黑、达堂和阿彼兰起来反抗他。他们不是灭亡了吗？让每个人各守其职。因为如果他教导的是谬误的道理，即使他是天使，也不要听从；但如果他教导真理，不要看他的生活，而要听他的话。你们有保禄用言语和行动教导你们正理。但你们说：\BibleQuote{他不周济穷人，他治理不善。}你们从哪里知道这些？在未了解之前不要责备。要畏惧那重大的交账。许多判断仅凭臆测。效法你们的主，祂说：\BibleQuote{我要下去，看看他们所行的，是否如我所闻，若不然，我也知道。}\BibleRef{创 18:21}但如果你们已经查问、了解、看见了，仍要等待法官，不要篡夺基督的职权。这是祂的事，不是你们的事。你们是卑微的仆人，不是主人。你们是羊，不要对牧人好奇，免得你们要为对他的指控交账。但你们说：他自己不实行，怎能教导我？不是他对你们说话。如果你们服从的是他，你们就没有赏报。是基督在训诫你们。我还说什么呢？你们甚至不应服从保禄，如果他说的是自己的话或属人的话，而应服从那有基督在他内说话的宗徒。让我们不要彼此论断行为，而是各人论断自己。省察自己的生命。\par

但你却说：\Quote{他应该比我好。}为什么？\Quote{因为他是司铎。}难道他在劳苦、危险、焦虑的争斗和患难上不比你更卓越吗？但如果他并不更好，你就因此毁掉自己吗？这是傲慢的话。他怎会不比你更好呢？你说：他偷窃，还犯亵渎圣物之罪！你怎么知道？你为什么把自己推下悬崖？假如你听说某人穿紫袍，你明知是事实，且能定他的罪，你却不肯去做，假装不知，不愿招惹不必要的危险。可是在这种情况下，你非但不退缩，反而无端将自己暴露于危险之中。不要以为这些话没有责任。听基督说的话：\BibleQuote{凡人们所说的每一句闲话，在审判之日，都要交账。}\BibleRef{玛 12:36} 你难道以为自己比别人好，而你却不像税吏那样叹息、捶胸、低头吗？\par

而纵使你更好，你却毁掉自己。沉默吧，这样你才不致失去更好。你若说出口，就抹杀了功德；你若这样想——我并非说你想——你若不想，就增添了许多。因为如果一个大罪人承认后\BibleQuote{回家成了义}，那么一个罪行较轻并自觉有罪的人，岂不更得赏报？省察你自己的生命：你不偷窃，但你贪婪、傲慢，犯了许多其他这类的事。我说这些并非为偷窃辩护——绝对不可！若真有人犯此罪，我深为惋惜，但我不相信。亵渎圣物之罪有多大，难以言表。但我宽恕你，因为我不愿我们的德行因指责他人而落空。有什么比税吏更坏的呢？他确实是税吏，犯了许多过犯，但因法利塞人只说了一句\BibleQuote{我不像这个税吏}，就毁了自己的一切功德。你说：我不像这个亵渎圣物的司铎。你岂不是也使一切落空了吗？\par

我不得不这样说，并在讲论中加以发挥，并非为我为他们挂虑，而是怕你们因自夸和论断别人而使自己的德行落空。请听保禄的劝诫：\BibleQuote{各人应当考验自己的行为，这样，他只在自身有所夸耀，而不在别人身上。}\BibleRef{迦 6:4}\par

如果你有一个伤口，请问，你去医生那里，你会阻止他给你的伤口敷药包扎，而好奇地询问医生自己是否有伤口吗？如果他真有伤口，你会在意吗？或者因为他有伤口，你就任由自己的伤口不治，说：医生应该身体健康，既然他不健康，我就让我的伤口不治了？难道对于受管辖的人来说，他的司铎是恶的，这能成为他的藉口吗？绝不。他将承受注定的惩罚，你也会遭遇你应得的报应。因为教师现在只是占据一个位置。因为经上记载：\BibleQuote{他们都要蒙天主教导。}\BibleRef{若 6:45}\BibleQuote{他们不必再教导人说：认识上主，因为不论大小，众人都要认识我。}\BibleRef{耶 31:34}那么，你会问，为什么他主持？为什么他被派来管理我们？我恳求你们，不要毁谤我们的教师，也不要对他们苛求，免得我们自招灾祸。让我们省察自己，我们就不会毁谤别人。让我们尊敬那一天，他光照了我们。人若有父亲，无论他有什么缺点，都会掩盖一切。因为经上说：\BibleQuote{不要因你父亲的羞辱而自夸，因为你父亲的羞辱并非你的荣耀；如果他心智衰退，你要容忍他。}\BibleRef{德 3:10}如果这话是对我们的生身父亲说的，那么对我们的神师更当如此。尊敬他，因为他每天为你服务，宣读圣经，为你整顿家庭，为你警醒，为你祈祷，为你恳求天主，为你呈上祈祷，他所有的敬礼都是为你。尊敬这一切，思想这些，并以虔诚的敬意接近他。不要说：他是恶的。那又如何？难道不是恶的人，就能凭自己将这些大恩惠赐予你吗？绝不。一切都按照你的信德运行。即使是义人，如果你不信，也不能使你受益；即使是不义的人，如果你忠信，也不能加害于你。天主在拯救祂的子民时，曾藉着牛为约柜行事。\BibleRef{撒上 6:12}难道司铎的善行或德行能赐予你如此大的恩惠吗？天主所赐的恩赐并非由司铎的德行所产生。一切都是恩宠。他的部分只是开口，而天主成就一切：司铎只是履行一个标记。试想若翰与耶稣之间相距何等遥远。听若翰说：\BibleQuote{我需要受祢的洗}\BibleRef{玛 3:14}，以及\BibleQuote{我连解祂的鞋带也不配。}\BibleRef{若 1:27}然而，尽管有这差距，圣神却降临了。这是若翰所没有的。因为经上说：\BibleQuote{从他的满盈中，我们都领受了。}\BibleRef{若 1:16}但是，圣神并未在祂受洗之前降临。然而，也并非若翰使祂降临。那么，为何这样发生呢？为叫你们知道，司铎履行一个标记。人与人之间的差异再大，也比不上若翰与耶稣的差异，然而圣神却与他同在并降临，使我们知道，是天主成就一切，一切都是天主的作为。我即将说出的话可能显得奇怪，但你们不要惊讶或惊惶。祭品是相同的，无论是普通人，还是保禄或伯多禄奉献。这与基督赐给祂的门徒的，以及现在司铎所祭献的是相同的。这丝毫不逊于那祭献，因为圣化此祭品的不是人，而是那圣化前者者也圣化了后者。因为正如天主所说的话与现在司铎所说的话相同，同样，祭品相同，圣洗也相同，都是祂所赐的。因此，一切都是出于信德。圣神立刻降临在科尔乃略身上，因为他先前已尽了自己的本分，献上了他的信德。这正是祂的身体，与那祭献一样。认为一个逊于另一个的人，不知道基督现今仍在、现今仍在运作。因此，知道这些事，我们并非无故说这话，而是为使你们的心灵归正，并使你们将来更加安全，要牢记所说的话。因为如果我们总是听道而不行道，就不会从这些话中获益。因此，让我们殷勤地听从所讲的话。让我们把它们铭刻在心。让它们永远印在我们的良心上，让我们不断光荣圣父、圣子及圣神。\par

\SourceRecord{Translated by Philip Schaff.；Nicene and Post-Nicene Fathers, First Series；Vol. 13.；Edited by Philip Schaff.；Buffalo, NY: Christian Literature Publishing Co.,；1889.；<http://www.newadvent.org/fathers/230702.htm>.}

% structure: p0001 source=h1 target=section reason=page-title

\section{弟茂德后书讲道第三篇}

\BibleRef{弟后 1:13-18}\par

\BibleQuote{你要持守从我这里听来的健全言辞的规范，并要持守在基督耶稣内的信德与爱德。你所受託的这美好事物，要靠那住在我们内的圣神来持守。你知道，凡在亚细亚的人都离弃了我，其中有菲革罗和赫摩革乃。愿主赐仁慈给敖尼息佛洛的家，因为他屡次使我精神焕发，不以我的锁链为耻；相反，当他到了罗马，他殷勤地寻找我，并找到了我。愿主赐他在那日子从主那里获得仁慈。他在厄弗所为我服务了多少事，你知道得最清楚。}\par

\Person{保禄}不仅以书信教导他的门徒尽本分，先前也以言语教导，这他在许多其他段落中表明，如他说：\BibleQuote{或是藉言语，或是藉我们的书信}\BibleRef{得后 2:15}，尤其是这里。因此，我们不要以为任何关于道理的教导是不完全的。因为他交付给他许多没有写下来的事。因此他提醒他，当他说：\BibleQuote{你要持守从我这里听来的健全言辞的规范。}如同艺术家的方式，我已将德行的形象印在你身上，在你灵魂中固定了一种规则、模型和一切中悦天主之事的轮廓。那么，你应持守这些，无论你默想信德或爱德，或健全心智的任何事，都从此形成你的观念。不必求助于他人作榜样，因为一切已存放在你内。\par

\BibleQuote{你要持守所受託的美好事物，}——如何？——\BibleQuote{藉着那住在我们内的圣神。}因为人的灵魂，当领受了如此伟大的事物时，本身不足以保证持守它们。为什么？因为有许多盗贼，浓厚的黑暗，魔鬼仍近在咫尺图谋害我们；我们不知道何时、何机会他会攻击我们。那么，他意思是，我们如何能足以持守它们？\BibleQuote{藉着圣神}；也就是说，如果我们有圣神与我们同在，如果我们不驱逐恩宠，祂会站在我们旁。因为，\BibleQuote{除非主建造房屋，建屋的人徒然劳苦；除非主守护城池，守城的人徒然警醒。}\BibleRef{咏 126:1}这就是我们的墙，我们的城堡，我们的避难所。因此，如果祂住在我们内，且祂自己就是我们的守卫，何需诫命？为使我们可以紧握祂、持守祂，而不以我们的恶行驱逐祂。\par

然后他描述自己的考验和诱惑，不是为了沮丧他的门徒，而是为了振作他，使他若将来遭遇同样的事，当回顾并记起他的老师所经历的事时，不致觉得奇怪。那么他说什么？既然弟茂德很可能被捕、被遗弃、得不到友善的关注、影响或援助，甚至被他的朋友和信徒们本身遗弃，听听他说：\BibleQuote{你知道，凡在亚细亚的人都离弃了我。}看来当时在\Place{罗马}有许多来自\Place{亚细亚}地区的人。\BibleQuote{但是没有人站在我一边，}他说，没有人承认我，所有人都疏远我。注意他灵魂的智慧。他只提他们的行为，并没有咒诅他们，反而赞扬那善待他的人，并为他祈求无数祝福，没有对他们发任何咒诅。\BibleQuote{其中有菲革罗和赫摩革乃。愿主赐仁慈给敖尼息佛洛的家，因为他屡次使我精神焕发，不以我的锁链为耻；相反，当他到了罗马，他殷勤地寻找我，并找到了我。}请注意他处处提到羞耻，而不提危险，以免弟茂德害怕。然而，那是充满危险的事。因为他与\Person{尼禄}的一个囚犯交朋友而得罪了\Person{尼禄}。但当他在\Place{罗马}时，他说，他不仅不回避与我交往，反而\BibleQuote{殷勤地寻找我，并找到了我。}\par

\BibleQuote{愿主赐他在那日子从主那里获得仁慈。他在厄弗所为我服务了多少事，你知道得最清楚。}\par

信徒应当如此。恐惧、威胁或耻辱，都不应阻止他们互相帮助，站在他们旁边，像在战争中援助他们。因为他们帮助那些在危险中的人，并不如他们自己得益，藉着为他们服务，使他们自己成为那些人的冠冕的分享者。例如，若有一个献身于天主的人遭受磨难和痛苦，并以极大的坚忍坚持战斗，而你尚未被卷入这战斗？你若愿意，不进入赛道，也能成为为他存留的冠冕的分享者，藉着站在他旁边，预备他的心思，鼓励和激发他。因此\Person{保禄}在其他地方说：\BibleQuote{你们与我同受苦难，做得很好。因为甚至在得撒洛尼，你们一次再次地送东西供应我的急需。}\BibleRef{斐 4:14}那些远离的人怎能与不在他们身边的人同受苦难呢？如何？他说：\BibleQuote{你们一次再次地送东西供应我的急需。}他又论到\Person{厄帕夫罗狄托}说：\BibleQuote{因为他几乎要死，不顾性命，为补足你们为我服务上的欠缺。}\BibleRef{斐 2:30}因为在君王的服役中，不仅那些作战的人，就连那些守护辎重的人，也分享荣耀；不但如此，还常常享有战利品的等分，尽管他们没有手染鲜血，没有列阵，甚至没有看见敌人的队伍；在这些战斗中也是如此。因为那当作战者饥饿时供应他、站在他旁边、用言语鼓励他并给予他一切服务的人，并不亚于作战者。\par

因为不要以为保禄这位战士、这位不可抗拒且不可战胜的战士，是许多人之中的某一位，他若没有受到许多安慰和鼓励，也许就不会站立得住，不会奋战。因此，那些不在竞赛中的人，或许可以成为正在竞赛者的得胜之因，并分享为胜者所保留的华冠。而且，如果那与活着的人相通的人被认为配与那些奋战者同享赏报，这有什么稀奇呢？因为在死后甚至也可以与已故者相通，与那些安眠、已得华冠、一无所缺的人相通。因为请听保禄的话说：\Quote{分享对圣人们的纪念}。这又如何做到呢？当你钦佩一个人，当你做了他因此而得华冠的那些事之一时，你显然与他的辛劳和他的华冠有份。\par

\Quote{愿主赐给他，使他在那一天得蒙主的仁慈}。他说，他曾怜悯了我，因此在那可怕的日子，当我们极需大量仁慈的时候，他也必将获得同样的回报。\Quote{愿主赐给他从主那里得蒙仁慈}。难道有两位主吗？决不是。而是\Quote{为我们只有一个主耶稣基督，并只有一个天主}。\BibleRef{格前 8:6}在此，那些受马西翁异端感染的人攻击这句话；但让他们知道，这种说法方式在圣经中并不罕见；如经上说：\Quote{上主对我主说}\BibleRef{咏 109:1}；又说：\Quote{我对上主说，你是我的主}\BibleRef{咏 15:2}；以及\Quote{上主从天上主那里降下火来}\BibleRef{创 19:24}。这表示诸位的本体相同，并非本性的区分。因为我们并非要理解为有两个彼此不同的本性，而是有两位，各自具有同一本体。\par

也要注意，他说：\Quote{愿主赐给他仁慈}。因为正如他自己从敖乃息佛罗那里得到了仁慈，同样他也愿他从天主那里得到仁慈。道德教训：如果敖乃息佛罗冒生命危险，是因仁慈而得救，那么我们就更因同一仁慈而得救了。因为那账目实在是可怕的，可怕的，并且极需要天主的仁爱，使我们不致听到那可怕的判决：\Quote{离开我……我从来不认识你们，你们这些作恶的人}\BibleRef{玛 7:23}；或那可怕的话：\Quote{可咒骂的，离开我，到那给魔鬼和他的使者预备的永火里去吧}\BibleRef{玛 25:40}；使我们不致听到：\Quote{在我们与你们之间，有巨大的深渊固定着}\BibleRef{路 16:16}；使我们不致听到那充满恐怖的声音：\Quote{把他丢在外面的黑暗中}；使我们不致听到那些充满恐惧的话：\Quote{你这又恶又懒的仆人}\BibleRef{玛 22:13}。因为那审判席实在是可畏而恐怖的。然而天主是仁慈而宽厚的。祂被称为\Quote{仁慈的天主和安慰的天主}\BibleRef{格后 1:3}；是至善的，无人像祂那样善，是仁慈、温柔、充满怜悯的，祂\Quote{不愿意罪人死亡，而愿意他悔改并生活}\BibleRef{则 18:24}。那么，那充满痛苦和苦楚的日子从何而来？一条火河在祂面前奔流。我们行为的书卷被打开了。那日子本身像火炉一样燃烧，天使们飞翔环绕，许多火窑已备好。那么祂怎么是仁慈、宽厚、充满对人的慈爱呢？正是在这一点上祂是仁慈的，并藉这些事显示祂慈爱的伟大。因为祂向我们提出这些恐怖之事，为的是使我们被它们约束，从而被唤醒去渴望天国。\par

并且要注意，除了赞扬敖乃息佛罗之外，他特别提及他的仁慈：\Quote{他屡次使我精神振作}；如同一个因炎热而疲惫的摔跤手，他在患难中使他精神振作、坚强起来。\newline{}\Quote{并且他在厄弗所怎样多方面地服事了我，你是十分知道的。}不只在厄弗所，而且在这里他也使我精神振作。因为一个警醒并醒觉于行善的人，应当有这样的行为，不是一次、两次、三次，而是终身如此。正如我们的身体不是一次吃饱，就可以供应一生的滋养，而是需要每日的饮食；同样，在此，虔敬也需要每日藉善工来维持。因为我们自己也需要极大的仁慈。正是由于我们的罪，天主、人类的挚友才行这一切，并非祂自己需要它们，而是祂为我们行了一切。因此，祂才将它们全都启示出来，并使我们知道它们，而且不仅告诉我们，还藉祂所行的事给予我们确证。虽然祂单凭祂的话就足以可信，但为使无人以为这是夸大其词或仅仅是一种威胁，我们还藉祂的作为得到进一步的保证。怎样？藉祂公开和私下所施的惩罚。为使你从这些实例中学习，祂有时惩罚了法郎，有时使洪水漫地并带来彻底的毁灭，又有时降下火雨：甚至如今，我们在许多事例中看到恶人遭受报应和惩罚，这些事是地狱的预像。\par

免得我们昏睡懈怠，忘记祂的话语，祂藉行动唤醒我们的心智；甚至在此世就向我们展示法庭、审判座和审讯。难道在人世间对正义竟如此重视，而天主——这些事物本就是祂所设立的——却毫不顾念吗？这又怎能令人信服？在家中，在市场，都有法庭。主人每日审判他的奴仆，追究他们的过失，惩罚一些，宽恕另一些。在乡间，农夫与他的妻子每日彼此论理。在船上，船长是法官；在营中，将军审判士兵；处处可见审判的程序。在各行各业中，师傅评判学徒。总之，所有人，公开地或私下地，彼此互为法官。任何事上，对正义的考量都未被忽视；各地的人都要为自己的行为交账。难道在此世对正义的追究如此遍及城市、家庭和个人，而在那里，在\BibleQuote{天主的右手充满正义}\BibleRef{咏 47:10}，在\BibleQuote{祂的正义有如天主的山岳}\BibleRef{咏 35:6}之处，反而对正义毫不顾念吗？\par

那么，为何天主——\BibleQuote{正义的审判者，刚强而忍耐}\BibleRef{咏 7:11}——这样容忍世人，而不立即施罚呢？在此你便知其原因：祂是宽仁的，藉此引导你悔改。但你若继续犯罪，那就是\BibleQuote{依据你的顽固和不悔改的心，为自己积蓄忿怒}\BibleRef{罗 2:5}。如果祂是正义的，祂就按应得的分施报，不忽视那些受冤屈的人，反而为他们报仇。因为这是正义者当行之事。如果祂是全能的，祂就在死后、在复活时予以报应，因为这是全能者的事。如果因祂的宽仁，祂容忍人类，我们就不应困惑，也不应问：为何祂不在此世施行惩罚？因为如果祂每日追讨我们的过犯，整个人类早已被扫除殆尽，因为确实没有一天是毫无罪污的，我们总在大事或小事上犯罪；以至于若无祂极大的宽仁和祂的良善——这良善赐予我们更长的悔改时间，好让我们除去以往的过犯——我们中没有一人能活到二十岁。\par

因此，让每个人凭着正直的良心，审视自己所行的一切，并将全部生活摆在眼前，反省他是否应受数不胜数的惩戒和刑罚？当他愤愤不平，认为某人犯了许多恶行却逍遥法外时，让他思量自己的过失，他的愤慨就会止息。因为那些罪行之所以显得重大，是由于发生在重大而昭彰的事上；但如果他省察自己的过失，或许会发现它们更为众多。因为偷盗和欺诈，无论对象是黄金还是白银，都是同一回事，因为两者出自同样的心念。一个愿偷少许的人，若有机会，也不会拒绝偷取许多；他没有偷，并非出于自愿，而是出于偶然的情况。一个穷人，若抢劫一个更穷的人，若能够的话，也必会毫不犹豫地抢劫富人。他的克制出于能力不足，而非选择。你说，某人是个统治者，抢掠治下之人的财产。但请问，你不是也在偷盗吗？因为不要告诉我，他偷的是塔冷通，而你偷的是几文钱。在施舍时，有人投下金子，而寡妇投了两文钱，但她所献的并不少于他们。为什么？因为看的是心意，而非礼物的多寡。那么，在施舍的事上，你愿意天主这样判断，并期望因你的贫穷，因投下两文钱而获得与投下许多塔冷通金子者同等的赏报？同样的规则不也适用于不义的行为吗？这又如何说得通呢？正如那投了两文钱的寡妇因她的善意而被视为与最慷慨的施舍者相等，同样，你偷了两文钱，便与那些更强大的强盗同样有罪。不仅如此，若容许我说出奇怪的话，你比他们是更恶劣的强盗。因为一个人，无论与国王的妻子、穷人的妻子或是奴隶的妻子犯奸淫，都是同等的奸夫：因为罪行不是按当事人的地位来评判，而是按犯罪者意志的邪恶程度；这件事也是如此。不，我甚至认为，与地位较低者犯此罪的人，或许比与王后本人勾搭的人更应受罚。因为后者可推说财富、美貌及其他诱惑，而前者却全无这些。因此，前者是更恶劣的奸夫。再者，在我看来，一个用劣酒达到酩酊大醉的人，是更顽固的酒徒；同样，一个不轻视小偷小摸的人，是更恶劣的诈骗者。因为犯大抢掠的人，或许不屑于小偷小摸；而偷小东西的人，绝不会停止更大的偷窃，因此，在两者中他是更大的贼。因为一个不轻视银子的人，怎能轻视金子呢？因此，当我们指控统治者时，让我们数算自己的过失，便会发现自己比他们更倾向于行不义和抢夺；除非我们按行为而非按心念的意向来判断是非，而事实上我们本应如此判断。如果一个人被认定偷了穷人的财物，另一个人偷了富人的财物，难道他们不是受同样的惩罚吗？一个人无论杀死的是贫穷丑陋的人，还是富有英俊的人，不都是同等的杀人犯吗？因此，当我们说某人强夺了他人的土地时，让我们反省自己的过失，这样就不会定别人的罪，反而会叹服天主的宽仁。我们不会因刑罚未降在他们身上而愤慨，反而会更加谨慎，不让自己行恶。因为当我们意识到自己也要受同样的惩罚时，就不再感到不满，并会停止犯罪，从而获得将来的恩惠，因我们的主耶稣基督的恩宠和慈爱，愿光荣归于祂与圣父，等等。\par

\SourceRecord{Translated by Philip Schaff.；Nicene and Post-Nicene Fathers, First Series；Vol. 13.；Edited by Philip Schaff.；Buffalo, NY: Christian Literature Publishing Co.,；1889.；<http://www.newadvent.org/fathers/230703.htm>.}

% structure: p0001 source=h1 target=section reason=page-title

\section{弟茂德后书第四篇讲道}

\BibleRef{弟后 2:1-7}\par

\BibleQuote{因此，我的儿，你要在\Person{基督耶稣}的恩宠上坚强。你从我在许多证人面前所听的，应托付给忠信的人，使他们也能教导别人。因此，你要如同\Person{耶稣基督}的好士兵一样，受苦难。没有当兵的而分心于世俗的事务，而能取悦于招募他的元帅。若有人参加比武，除非依法比武，否则得不到花冠。勤劳的农夫应先得到果实。你要反思我所说的；愿主赐你在一切事上明悟。}\par

年轻的海员在海上，若船长在船难中得以保全，会激发极大的信心。因为他不会认为他遭遇风暴是由于自己缺乏经验，而是出于事物的本性；这一点对他的思想影响颇大。在战争中，上尉看见将军受伤而又痊愈，也大受鼓励。同样，宗徒曾遭受极大的苦难，却未被最极端的苦难削弱，这给信友带来某种安慰。若不是这样，他也不会讲述自己的苦难。因为当\Person{弟茂德}听到那位拥有如此大能、征服了全世界的人，成了囚犯，遭受苦难，却不焦躁，也不因朋友的背弃而不满；那么，他若自己也遭受同样的苦难，就不会认为这是出于人性的软弱，也不会因为自己是门徒、比\Person{保禄}低微（因为他的老师也受过类似的苦）而认为这出于某种缘由，而是认为这一切都是出于事物的本性。因为\Person{保禄}本人就这样做了，讲述了自己的遭遇，为的是坚强\Person{弟茂德}，更新他的勇气。他表明他正是为此提及自己的考验和磨难，因为他接着说：\BibleQuote{因此，我的儿，你要在\Person{基督耶稣}的恩宠上坚强。}你说什么？你用恐惧动摇了我们，你告诉我们你身陷锁链、处于磨难，众人都离弃了你；然后，好像你什么也没受过、也未被任何人遗弃似的，你又加上：\BibleQuote{因此，我的儿，你要坚强}？这也是理所当然的，因为这些事情为的是坚强你，而非坚强他。因为如果我这\Person{保禄}都能忍受这些事，你更应当承受它们。如果老师如此，门徒更当如此。而且他带着深厚感情引出这一劝诫，称他\Quote{儿子}，不仅这样，还说\Quote{我的儿子}。他的意思是：如果你是一个儿子，就要效法你的父亲。如果你是一个儿子，就要因我所说过的话而坚强；或者更确切地说，要坚强，不仅出于我对你说的，而是\Quote{出于天主}。他说：\BibleQuote{要在\Person{基督耶稣}的恩宠上坚强}；也就是说，\Quote{藉着\Person{基督}的恩宠}。那就是要坚定地站立。你知道这场战斗。因为他在别处说：\BibleQuote{我们不是与血肉搏斗。}\BibleRef{弗 6:12}他说这话并非要打击他们，而是激励他们。因此他的意思是：要清醒，要警醒，有主的恩宠与你合作，帮助你在竞赛中，你也要以极大的热忱和决心贡献你的部分。\BibleQuote{你从我在许多证人面前所听的，应托付给忠信的人}；托付给\Quote{忠信}的人，不是给好问的人，也不是给好推理的人，而是\Quote{忠信}的人。怎样的忠信？就是不背叛他们所应宣讲的福音的人。\BibleQuote{你所听的}，不是你探究出来的。因为\BibleQuote{信德来自听闻，听闻来自天主的圣言。}\BibleRef{罗 10:17}但为什么说\BibleQuote{在许多证人面前}？仿佛说：你并非私下或独自听闻，而是在众人面前，以完全的公开方式。他也不是说\Quote{说出}，而是说\Quote{托付}，如同将委托的宝藏安全存放。他又从高处和低处警戒他的门徒。但他不仅说\BibleQuote{托付给忠信的人}；因为一个人若是忠信，却不能将他的道理传授给别人，这有什么益处呢？他固然不背叛信德，但不能使别人也成为忠信的吗？因此教师应具备两种品质：既忠信，又善于教导；所以他接着说：\BibleQuote{使他们也能教导别人。}\par

\BibleQuote{因此，你要如同\Person{耶稣基督}的好士兵一样，忍受苦难。}哦，这是何等大的尊严，成为\Person{耶稣基督}的士兵！请看地上的君王，在他们麾下服役被视为多么大的荣誉。因此，如果君王的士兵应当忍受苦难，那么不忍受苦难就不是任何士兵的本分。所以，你若忍受苦难，不应抱怨，因为那是士兵的本分；但若不忍受苦难，才应抱怨。\par

\BibleQuote{没有当兵的而分心于世俗的事务，而能取悦于招募他的元帅。若有人参加比武，除非依法比武，否则得不到花冠。}\par

这些话固然是对\Person{弟茂德}说的，但通过他，也是对每一位教师和门徒说的。因此，凡担任主教职位的，谁也不应不屑于听这些话，而若不去行，就应感到羞愧。\Quote{若有人参加比武，}他说，\Quote{除非依法比武，否则得不到花冠。}所谓\Quote{依法}是什么意思？他进入竞技场、受过膏油、甚至参与比赛还不够，除非他遵守所有训练规则，包括饮食、节制与清醒，以及角力学校的一切规则；简言之，除非他经历一切角力者应做之事，否则他得不到花冠。请观察\Person{保禄}的智慧。他提到角力者和士兵：前者是为使\Person{弟茂德}准备面对杀戮和流血，后者是为忍耐，好使他能刚毅地承受一切，并经常操练。\par

\BibleQuote{勤劳的农夫应先得到果实。}\par

他先是以自己作为教师的榜样说话。如今他更从一些普通的事例说话，如同角力者和兵士，并且在他们身上，他把奖赏摆在他面前。第一，为使他能取悦于那选他当兵的人；第二，为使他能得冠冕；现在他提出第三个更特别适合他自己的例子。因为兵士和角力者的例子对应那些受管辖的人，而农夫的例子则是对应教师。（要努力）不仅作为兵士或角力者，也要作农夫。农夫不只照顾自己，也照顾大地的果实。这就是说，农夫从他的劳动中享得不少报酬。\par

这里他既表明，对天主而言一无所缺，又表明教导是有赏报的，他以一个普通例子来显明这事。他说，正如农夫不是无利益地劳动，而是比别人先享受自己劳苦的果实，同样教师也应当如此：他或是指这个意思，或是指当给教师尊荣，但这不太一致。因为他为何不单说\Quote{农夫}，而说\Quote{劳苦的呢}？不只是那工作的，更是那劳累至疲惫的？于此，关于奖赏的延迟，为使无人急躁，他说，你们已收割果实，或说劳苦本身就有赏报。因此，当他将兵士、角力者和农夫的例子，以及所有比喻性的例子摆在他面前时，他说：\Quote{没有人}能得冠冕，\Quote{除非他合法地竞赛}。且他观察到\Quote{劳苦的农夫必须首先享受果实}之后，他接着说：\par

\Quote{你要思想我所说的话，愿主在一切事上赐给你明悟。}\par

正是为此，他用谚语和比喻说了这些话。然后又为显示他的慈爱心肠，他不停地为他祈祷，如同为自己的儿子担忧，并说：\par

8, 9节。\Quote{你要记得耶稣基督，出自达味的后裔，从死中复活了，是按照我的福音。我因此受苦，如作恶的，甚至被捆绑。}\BibleRef{弟后 2:8-9}\par

\Quote{你要记得耶稣基督，出自达味的后裔，从死中复活了。按照我的福音。} 因此他在各书信中处处说\Quote{按照我的福音}，或因他们必须相信他，或因有些人传讲\Quote{别的福音}。\BibleRef{迦 1:6}\par

\Quote{我因此而受苦，如同作恶的，甚至被捆绑。}他又从自身引入安慰和鼓励，用这两件事预备听者的心：第一，要知道他忍受艰苦；第二，他不是徒然如此，而是为有益的目的，因为在此他将获益，否则甚至受害。因为如果你能表明一位教师遭受艰难，却无任何有益目的，那有什么益处呢？但若是为任何益处，为被教导者的利益，那就值得赞美。\par

\Quote{但天主的圣言并不被捆绑。} 这就是说，如果我们属于这世界的兵士，从事属地的战争，那么捆住我们双手的锁链就起作用。但现在天主使我们成为无人能制服的了。因为我们的手被捆绑，但我们的舌不受捆绑，因为除了怯懦和不信，没有什么能捆绑舌头；而在没有这些的地方，即使你把锁链加在我们身上，福音的宣讲也不被捆绑。如果你捆绑一个农夫，你就阻止了他撒种，因为他用手撒种；但如果你捆绑一位教师，你并不阻碍圣言，因为圣言是用舌头撒种，而不是用手。因此我们的圣言不受锁链的束缚。我们虽被捆绑，圣言却是自由的，并运行它的路程。如何呢？因为我们虽被捆绑，看，我们仍在宣讲。这是对那自由之人的鼓励。因为我们被捆绑的尚且宣讲，何况你们这些得释放的，更应该这样做。你们已听说我因这些事受苦，如同作恶的。不要沮丧。因为这是一个大奇事：被捆绑的我却做那自由之人的工作，被捆绑的我战胜一切，被捆绑的我压倒那捆绑我的人。因为这是天主的圣言，不是我们的。人间的锁链不能捆绑天主的圣言。\Quote{我为蒙选的人忍受这一切。}\par

10节。\Quote{因此我为万物忍耐，}他说，\Quote{为蒙选者的缘故，为使他们也获得那在基督耶稣内的救恩，连同永远的荣耀。}\BibleRef{弟后 2:10}\par

请看另一个激励。他说：我忍受这些，不是为了我自己，而是为了他人的救恩。我本可以生活在危险之外；如果我只顾自己的利益，就无需遭受这些。那么我为什么受这些苦？为了别人的好处，使别人能获得永生。那么你给自己许诺了什么？他没有简单地说为了这些特定的人，而是\Quote{为了被选者的缘故}。如果天主选择了他们，我们就应当为他们承受一切。\Quote{使他们也能获得救恩}。他说\Quote{他们也能}，意指像我们一样。因为天主也选择了我们；正如天主为我们受苦，我们也该为他们受苦。因此这是回报，而非恩惠。在天主方面，是恩宠，因祂未受任何先前的恩惠而善待了我们；但在我们方面，是回报，我们先前从天主领受了恩惠，就为这些人受苦，\Quote{使他们获得救恩}。你说什么？什么救恩？你连自己的救恩都不是作者，反倒毁坏自己，你能成为他人救恩的作者吗？当然不能，因此他加上\Quote{在基督耶稣内的救恩}；那是真正的救恩，\Quote{带着永久的荣耀}。现世的事是痛苦的，但只是在地上。现世的事是可耻的，但只是暂时的。它们充满苦涩和痛苦，但只持续今天和明天。\par

善事的性质并非如此，它们是永恒的，它们在天上。那是真正的荣耀，这是耻辱。\par

道德教训。请看，我恳求你们，亲爱的，那地上的不是荣耀，真正的荣耀在天上。但若有人要得荣耀，就让他受羞辱。若有人要得安息，就让他受苦。若有人要永远显赫，享受快乐，就让他轻看暂时的事物。并且羞辱就是荣耀，荣耀就是羞辱，现在让我们在能力范围内将这一点摆在我们面前，好看到什么是真正的荣耀。在地上受荣耀是不可能的；你若要受荣耀，必须通过羞辱。让我们以两个人的例子来证明这一点：尼禄和保禄。一个拥有这世界的荣耀，另一个拥有羞辱。怎么样的呢？第一个是个暴君，获得了巨大的成功，树立了许多凯旋柱，财富不断涌来，到处是庞大的军队；他拥有世界的大部分和帝国都城，服从他的统治，整个元老院都向他屈膝，他的宫殿也以辉煌的景象前进。当他必须武装时，他身着金子和宝石出行。当他安坐和平时，他穿着紫袍。他被众多的卫兵和侍从包围。他被称为陆海之君、皇帝、奥古斯都、凯撒、国王，以及其他诸如此类暗示奉承和谄媚的响亮名号；没有一样有助于荣耀的东西是缺乏的。甚至智者、权贵和君主都向他战栗。因为除了这一切，他还被认为是一个残忍和暴力的人。他甚至希望被当作神，他蔑视所有的偶像，也蔑视那凌驾万有的天主。他被当作神崇拜。还有比这更大的荣耀吗？或者不如说，还有比这更大的羞辱吗？因为——我不知道怎么回事——我的舌头被真理的力量带走，在审判之前就下了判决。同时，让我们根据大众、不信者以及奉承的评价来审视这件事。\par

在共同的荣耀观念中，还有什么比被视为神更大的呢？任何人类竟如此疯狂，这确实是极大的耻辱，但现在让我们按照大众的意见来考虑这件事。那么，凡有助于人类荣耀的东西，他没有一样是缺乏的，但他被所有人当作神崇拜。现在，与他对立，让我们考虑保禄。他是基里基雅人，罗马和基里基雅之间的区别，所有人都知道。他是一个制帐幕者，一个穷人，不精通外人的智慧，只懂得希伯来语，一种被所有人轻视的语言，尤其是被意大利人。因为他们并不那么轻视蛮族语、希腊语或任何其他语言，而是轻视叙利亚语，而叙利亚语与希伯来语相近。这并不奇怪，因为如果他们轻视那如此卓越和优美的希腊语，那么就更轻视希伯来语了。他是这样一个经常挨饿的人，常常不进食就上床睡觉，一个赤身露体、没有衣服可穿的人；正如他论自己所说的：\Quote{寒冷与赤身露体}。\BibleRef{格后 11:27}不仅如此，而且他被尼禄本人下令投进监狱，与强盗、骗子、盗墓者、杀人犯关在一起，并且正如他自己所说的，他作为犯人被鞭打。那么谁更杰出呢？一个人的名字，大部分人从未听说过。另一个人每日被希腊人、蛮族人、斯基泰人以及居住在地极的人称颂。\par

但我们暂且不论现今的情况，单看当时：谁更显赫，谁更光荣？是那被锁链捆绑、从狱中被拖出的囚犯，还是那身穿紫袍、从王宫中走出的君王？无疑是那囚犯。因为那统率军队、身穿紫袍的君王，并不能随心所欲；而那像罪犯一样、衣着卑微的囚犯，却能更有权威地行事。何以见得？前者说：\Quote{不可传布天主的圣言。}后者说：\Quote{我不能不说；\InnerQuote{天主的圣言决不会被束缚。}}于是，那位奇里基雅人，那囚犯，那贫穷的帐棚匠，忍饥挨饿，却藐视那富有的罗马人、皇帝、统治万有的人，他使成千上万的人致富；然而他虽有全部军队，却毫无用处。那么，谁更显赫？谁更可敬？是那戴着锁链却得胜的人，还是那穿着紫袍却败亡的人？是那站在下方却击倒对方的人，还是那坐在上方却被击倒的人？是那发令却受鄙视的人，还是那听命却不理会命令的人？是那独自一人却得胜的人，还是那率领大军却战败的人？那君王最终如此失败，以致他的囚犯战胜了他。那么请告诉我，你愿意站在哪一边？不要看后来发生的事，只看他们当时的状态。你愿意站在尼禄一边，还是保禄一边？我不是按信德来估计，因为那显而易见；而是按光荣、尊敬和优越来估计。凡有正确理解的人都会说：保禄一边。因为得胜比失败更光荣，所以他更光荣。但这一点还不算重要：他得胜了，而且是在如此卑微的状态下，战胜了如此尊贵的人。因为我要说，并且不会停止重复：虽然戴著锁链，他却击倒了那戴着王冠的人。\par

这就是基督的能力。那锁链胜过了君王的冠冕，这装束显得比那冠冕更灿烂。他穿着囚犯的褴褛衣衫，却使所有人的目光都转向他身上的锁链，而不是那紫袍。他站在地上，被捆绑着，弯腰屈膝，而众人离开了那驾驭金色战车的暴君，来注视他。他们这样做是合理的。因为看见国王骑着白马是平常事，但看见一个囚犯以国王对待可怜卑微奴隶的自信与国王交谈，却是奇异而罕见的景象。周围的人群都是国王的奴隶，但他们不钦佩自己的主人，却钦佩那超越主人的人。那人人畏惧战栗的人，竟被一个孤独的人践踏。请看，这些锁链的光彩何其大啊！\par

何必再提此后的事呢？一人的坟墓无处可寻；但另一人却躺在皇城之中，比任何君王都更辉煌，甚至在他得胜、树立凯旋标志的地方。若提到一人，即使在他的亲属中也是带着羞辱，因为据说他荒淫无度。但另一人的记忆却到处伴随着美名，不仅在我们中间，也在他的敌人中间。因为真理彰显时，甚至使敌人羞愧；即使他们不为他的信德而钦佩他，也会为他的勇气和豪迈的自由而钦佩他。一人被众口传颂，如同戴冠者；另一人则被责骂和控告所淹没。那么，真正的光辉是什么呢？\par

但我不过是在称赞狮子的爪牙，而我本应讲述他真正的荣耀。那荣耀是什么呢？就是天上的荣耀。他将如何穿着闪亮的衣袍，与天上君王一同来临！那时尼禄将如何站立，悲伤而沮丧！如果我说的话在你看来不可思议、荒谬可笑，那么你嘲笑那不应被嘲笑的事，才是可笑的。因为如果你不信未来，就当从过去的事中得确信。受冠的季节尚未到来，然而这位战士已经获得何等大的光荣！那么当分发赏报者来临之时，他将获得何等大的光荣！他曾身在异乡，\Quote{是客人，是旅居者}\BibleRef{希 11:13}，就这样受到敬仰；当他回到自己的地方时，他将会享有多少美善！如今\Quote{我们的生命与基督一同藏在天主内}\BibleRef{哥 3:3}；然而那已死的人，比活着的人更工作、更受尊敬。当我们的生命显现时，他将不会参与什么？不会获得什么？\par

因此，天主使他享有这些尊荣，并非因为他需要它们。因为若他在肉身时轻视世俗的荣耀，如今他已脱离肉身，更会轻视它。不仅如此，天主使他享有尊荣，也是为使那些不信将来之事的人，能因现在之事而信服。我说，当复活之时，\Person{保禄}将与天国之王一同来临，并享受无限的福乐。但不信者不会信服。那么，让他从现在之事而信服吧。这帐棚匠比君王更显赫、更受尊荣。\Place{罗马}没有一位皇帝曾享有如此大的尊荣。皇帝被驱逐，埋没无人知晓。帐棚匠却占据城中心，如同一位活着的君王。从这些事，你当相信将来之事。若他在这里——他曾受迫害和被放逐之地——享有如此大的尊荣，那么当他日后来临之时，又将如何呢？若他身为帐棚匠时如此显赫，那么当他来临与太阳争辉之时，又将如何呢？若他在如此卑微中战胜了如此辉煌，那么在他来临时，谁能不被胜过呢？我们岂能避免这结论？有谁不被这事实所感动：一个帐棚匠变得比最受尊荣的君王更尊贵？若在此世事情如此超乎自然，那么来世更是如此。人啊，你若不信将来之事，就当信现在之事。你若不信不可见之事，就当信可见之事；或者更确切地说，信可见之事，这样你便会信不可见之事。但你若不肯，我们便可适宜地与宗徒同说：\BibleQuote{我们从你们的血中是无罪的}\BibleRef{宗 20:26}，因为我们已向你们见证了一切，没有遗漏任何该说之事。因此，你们当责怪自己，并将地狱的惩罚归于自身。但是，我亲爱的孩子们，让我们效法\Person{保禄}，不仅在他的信德上，也在他的生活上，好使我们能获得天上的荣耀，并践踏此世的荣耀。愿任何现世之事都不吸引我们。让我们轻视可见之事，为能获得天上的事物，或者更确切地说，藉由这些获得那一切。但让我们首要的目标是获得那些事物，愿天主赐予我们众人堪当承受，藉着恩宠和慈爱，等等。\par

\SourceRecord{Translated by Philip Schaff.；Nicene and Post-Nicene Fathers, First Series；Vol. 13.；Edited by Philip Schaff.；Buffalo, NY: Christian Literature Publishing Co.,；1889.；<http://www.newadvent.org/fathers/230704.htm>.}

% structure: p0001 source=h1 target=section reason=page-title

\section{论弟茂德后书 讲道 5}

弟茂德后书 2:11-14 \BibleRef{弟后 2:11}\par

\Quote{这是忠信的话：如果我们与祂同死，也必与祂同活；如果我们受苦，也必与祂一同为王；如果我们否认祂，祂也必否认我们；如果我们不信，祂仍是信实的，因为祂不能否认自己。你当把这些事提醒他们，在主面前劝诫他们，不要为言辞争辩，这毫无益处，只能败坏听的人。} \BibleRef{弟后 2:11}\par

许多较弱的人放弃了信德的努力，不能忍受希望的延迟。他们寻求眼前的事物，并以此来判断未来。因此，当他们在世上的命运是死亡、折磨和锁链，而他却说他们将获得永生时，他们不会相信，反而会说：\Quote{你说什么？我活着就死，死了就活？你在世上什么也不应许，却在天上应许？你不赏赐小事，却要赏赐大事吗？} 为使无人如此争辩，他预先将证明这些事的证据置於无疑之地，并且给出某些记号。因为他说：\Quote{你记起}，他说，\Quote{耶稣基督从死者中复活了}；也就是死而复活。现在显示同一件事，他说：\Quote{这是\OriginalTerm{忠信的话}{faithful saying}}，即得到天上生命的人也将得到永生。那里\Quote{忠信}是从何而来的呢？因为他说：\Quote{如果我们与祂同死，也必与祂同活。} 试问，我们既有分于祂辛劳痛苦的事，难道不会分于有益的事吗？但即使是人也不会这样做；如果有人选择与他一同受苦受死，当那人得到安息时，也不会拒绝让他分享。但我们如何\OriginalTerm{与祂同死}{dead with Him}呢？他所指的死，既是在\OriginalTerm{圣洗池}{Laver}中的死，也是受苦中的死。因为他说：\Quote{身上常带着主耶稣的死} \BibleRef{格后 4:10}；又说：\Quote{我们藉着圣洗与祂同葬于死} \BibleRef{罗 6:4}；又说：\Quote{我们的\OriginalTerm{旧人}{old man}与祂同钉十字架}；又说：\Quote{我们已与祂同植于祂死亡的相似中} \BibleRef{罗 6:5}。但他这里也提到因试炼而死的死：尤其如此，因为他写这话时正受着试炼。这正是他所说的：\Quote{如果我们为祂的缘故受了死，难道不会为祂的缘故活着吗？这是无可置疑的。\InnerQuote{如果我们受苦，我们也将与祂一同为王}}，不是绝对地说我们将为王，而是\Quote{如果我们受苦}，表明死一次是不够的（这位有福的人自己天天死），而需要极大的忍耐；弟茂德尤其需要。因为不要说，他说道，你最初的受苦，而说你要继续受苦。\par

然后另一方面，他劝诫他，不是从善的，而是从恶的。因为如果恶人也分得同样的事物，这就不是安慰。如果忍受的人将与祂一同为王，而不忍受的人确实不与祂一同为王，但也不会遭受更糟的恶，虽然这可怕，却不足令多数人关切。因此他说到更可怕的事：如果我们否认祂，祂也必否认我们。所以不仅有善的报应，也有相反的。试想那在国度中被否认的人，很可能要受什么苦？\Quote{凡否认我的，我也必否认他。} \BibleRef{玛 10:33} 这报应并不相等，尽管看起来是这样表达的。因为否认祂的我们是人，而否认我们的祂是天主；天主与人之间有多大的距离，无须多言。\par

此外，我们伤害自己；我们不能伤害祂。为显示这一点，他补充说：\Quote{如果我们不信，祂仍是信实的：祂不能否认自己。} 也就是说，如果我们不信祂复活了，祂并不因此受损。祂是信实而不动摇的，无论我们是否这样说。那么，既然祂毫不因我们否认祂而受损，祂渴望我们的宣认无非是为了我们的益处。因为无论我们否认祂与否，祂始终如一。祂不能否认自己，即祂的自身。我们可能说祂不存在；尽管事实并非如此。祂的本性不容许祂不存在，即归于虚无。祂的存有永远常存，总是存在。所以，我们不要以为我们可以取悦祂或伤害祂。但唯恐有人认为弟茂德需要这劝告，他又补充道：\par

\Quote{你当把这些事提醒他们，在主面前劝诫他们，不要为言辞争辩，这毫无益处，只能败坏听的人。} 呼求天主为我们所说的作证是一件威慑人心的事，因为如果没有人敢蔑视被呼求的人证，何况呼求天主呢？例如，如果有人进入合同，或立遗嘱，选择呼召可信的证人，谁会把财物转给未列入的人？当然不会。即使他愿意，但害怕证人的可信度，也回避这样做。什么是\Quote{在主面前劝诫他们}？他呼求天主为所说和所做的作证。\par

\Quote{不要为言辞争辩，这毫无益处}；不仅如此，还有\Quote{败坏听的人}。不但无益，而且有害。\Quote{那么，你当把这些事提醒他们}，如果他们轻视你，天主会审判他们。但他为什么告诫他们不要为言辞争辩？他知道这是一件令人愉快的事，人的灵魂总是倾向于争辩和争论言辞。为防止此事，他不仅命令他们\Quote{不要为言辞争辩}，而且为使他的言论更令人警惕，他补充说\Quote{败坏听的人}。\par

第15节。\Quote{你当竭力在天主面前呈上自己为蒙悦纳的工人，无愧的，\OriginalTerm{正确分剖真理之言}{rightly dividing the word of truth}。} \BibleRef{弟后 2:15}\par

随处都有这\Quote{不以为耻}！为何他如此小心地防范弟茂德以羞愧？因为对许多人而言，以保禄本人——一个制帐棚的——为耻，以及以那宣讲为耻（因其教师们都死了）是自然的。因为基督曾被钉在十字架上，保禄自己将被斩首，伯多禄被倒钉在十字架上，而这些事他们竟受自那些狂妄可鄙的人。由于那样的人掌权，他说：\Quote{不要以为耻}；意思是，不要因敬畏而不敢做任何促进虔敬的事，即使需要为奴或遭受任何其他苦难。因为一个人如何成为蒙悦纳的呢？是藉着成为\Quote{一个无需羞愧的工人}。正如工人不以任何工作为耻，同样，在福音上劳苦的人也不应以任何事为耻。他当承受一切。\par

\Quote{正确地分解真理之言。}\par

他说得好。因为许多人扭曲它，以各种方式曲解它，还添加了许多东西。他没有说\Quote{指导它}，而是说\Quote{正确地分解}，即砍掉伪造的部分，以极大的猛烈打击它，并根除它。用圣神的剑，从你的宣讲中，如同从皮条上，割掉一切多余和无关的东西。\par

第16节：\Quote{并躲避凡俗的空谈。} \BibleRef{弟后 2:16}\par

因为他们不会停在那里。因为一旦引入任何新东西，它就会不断产生革新，而一个曾离开安全港湾的人的错误是无限的，永不停歇。\par

\Quote{因为他们必将趋于更不虔敬，}他说，\par

第17节：\Quote{他们的话如同毒疮蔓延。} \BibleRef{弟后 2:17}\par

那是一种无法约束的恶，无药可治，它摧毁整个结构。他指出教义的新奇是一种疾病，比疾病更糟。这里他暗示他们是不可救药的，而且他们的错误不是出于软弱，而是出于故意。\par

\Quote{其中有依默纳约和非肋托，}\par

第18节：\Quote{他们偏离了真理，说复活已过，就败坏了一些人的信仰。} \BibleRef{弟后 2:18}\par

他说得好：\Quote{他们必将趋于更不虔敬。}因为这事看来似乎是一种孤立的恶，但看看从中生出多少恶事！因为如果复活已经过去，我们不仅因被剥夺了那大荣耀而受损，而且因审判被除去，报应也被除去。因为如果复活已经过去，报应也已过去。因此，善人已收获了迫害和苦难，恶人未受惩罚，反而生活在极大的享乐中。还不如说没有复活，倒比说复活已经过去更好。\par

\Quote{并且败坏，}他说，\Quote{了一些人的信仰。}\par

\Quote{一些人的，}不是所有人的。因为如果没有复活，信仰就被颠覆了。我们的宣讲是空的，基督也没有复活；而如果他没有复活，他也就没有出生，也没有升天。请看这错误，虽然看似反对复活的道理，却引来了许多其他恶事。那么，有人会说，我们不该为那些被败坏的人做些什么吗？\par

第19节：\Quote{然而，}他说，\Quote{天主的坚固基础屹立不移，且有这个印号：\InnerQuote{主认识那些属于他的人。}并且，\InnerQuote{凡称呼主名的人，应当远离不义。}} \BibleRef{弟后 2:19}\par

他指出，甚至在他们被败坏之前，他们就不坚固。否则，他们就不会在第一次攻击时就倾覆，正如亚当在诫命之前是坚固的。因为那些固定的人不仅不受欺骗者的伤害，甚至受钦佩。\par

他称之为\Quote{坚固的}和\Quote{基础}；所以我们应当持守信仰；\Quote{有这个印号：\InnerQuote{主认识那些属于他的人。}}这是什么意思？他从申命记中取来；意思是：坚固的灵魂坚定而不动摇。但他们如何显明呢？是由于这些特征刻在他们的行为上，是由于他们被天主认识，而不与这世界一同灭亡，并且是由于他们远离不义。\par

\Quote{凡称呼主名的人，}他说，\Quote{应当远离不义。}\par

这些是这基础的标志。正如基础被显明为坚固，并且如同字母刻在石头上，使字母有意义。但这些字母是由行为显明的，\Quote{有}他说，\Quote{这个印号}固定其上：\Quote{凡称呼主名的人，应当远离不义。}所以，如果有人不义，他就不属于这基础。因此，这也属于那印号：不做不义的事。\par

论道德。因此，我们不要从自己身上除去那王室的印号和标记，以免我们成为那些没有受印的人，以免我们不够坚固，以免我们动摇不定，反而要属于那基础，不被摇摆不定。那标记那些属于天主的人的是：他们远离不义。因为一个人怎么能属于那位正义的天主呢，如果他行不义，如果他的行为与他敌对，如果他的恶行侮辱他？我们再次谈论反对不义的事，而我们又有许多敌对者。因为这种情感，如同一个暴君，占据了所有人的灵魂；而更糟的是，它不是出于必要或暴力，而是出于说服和温和的暗示，而人们还对他们的奴役心存感激。这确实是可悲的事；因为如果他们是被束缚而非出于爱，他们很快就会离开。那么，最苦的东西怎么会变成甜的呢？最甜的公义又怎么会变成苦的呢？这要归咎于我们的感官。因此有些人以为蜂蜜是苦的，却高兴地吃其他有害的东西。原因不在于事物的本性，而在于受害者的歪曲。灵魂的判断官能紊乱了。正如一架天平，如其横梁不稳，就会摇摆，不能准确显示放在其中的重量；同样，灵魂若缺少其思想的横梁——这横梁固定并坚固地铆接于天主的法律，被摇摆、被下拉，就不能正确地判断行为。\par

因为如果有人仔细审视，他会看出不义的巨大苦涩，不是对受害者，而是对行不义者，而且对他们比其他人更甚。我们不要谈论将来的事，只谈现在在此世的事。难道它没有争斗、审判、定罪、恶意、辱骂吗？还有什么比这些更苦涩？难道它没有仇恨、战争、控告吗？还有什么比这些更苦涩？难道它没有良心不断地鞭挞和啃噬我们吗？如果可能，我愿意把不义者的灵魂从身体里拉出来，你会看见它苍白颤抖，羞愧，藏头，焦虑恐惧，自我定罪。因为即使我们沉入罪恶的深渊，心灵的判断力并未被摧毁，而是保持不受贿赂。没有人行不义以为那是善，而是编造借口，用尽一切言语的诡计来推卸指控。但他无法摆脱良心的指控。在这里，言辞的动听、统治者的腐败、众多谄媚者，往往能使正义蒙上阴影，但在内心，良心没有这类东西，那里没有谄媚者，没有财富来贿赂审判者。因为判断的能力是天主自然植入我们内的，出自天主的不能被如此败坏。但不安的睡眠，纷至沓来的幻想，以及对罪过的频繁回忆，破坏了我们的安宁。例如，有人不义地剥夺了别人的房子？不仅被抢劫的人不幸，抢劫他的人也不幸。如果他相信未来的审判（如果真有人这样相信），他会极其叹息，处于痛苦中。但如果他不相信来世，他仍会因羞愧而脸红；或者更确切地说，无论是希腊人、犹太人还是异端者，没有人不怕将来的审判。\par

尽管他对来世不是哲人，但他害怕并颤抖于在此世可能遭遇的事，唯恐他的财产、子女、家庭或生命受到某种报应。因为天主降下许多这样的惩罚。既然复活的道理不足以使所有人恢复理智，祂甚至在此世提供许多祂公义审判的证据，并向世界展示。一个非法获利的人没有子女，另一个在战争中倒下，另一个身体残缺，另一个失去儿子。他思考这些事，他的想象萦绕在这些事上，他生活在持续的恐惧中。\par

难道你不知道不义者所受的苦？这些事中没有苦涩吗？即使没有这类事，难道不是所有人都谴责他，恨恶他，认为他比野兽更不理性，甚至连那些自身不义的人也是如此吗？因为如果他们谴责自己，就更会谴责别人，称他为贪婪、欺诈、有害之人。那么他能享受什么快乐呢？他只有更沉重的忧虑和焦虑来保全他的所得，并且更加焦虑和烦恼。因为一个人积聚的财富越多，他为自己储存的痛苦的警戒就越多。然后那些被他冤枉的人的诅咒、他们对他的指控呢？如果疾病临到他呢？因为一个人一旦生病，无论他多么无神论倾向，他不可能不担心这些事，不可能不思想，当他什么都不能做的时候。因为只要我们在此世，灵魂享受自身，不容忍痛苦的思想；但当它即将离开身体时，更大的恐惧束缚它，如同进入审判的门槛。即使是强盗，当他们在监狱里时，生活没有恐惧，但当他们被带到法庭的帷幕前时，他们因恐惧而瘫软。因为当死亡的恐惧紧迫时，像火一样烧尽其他一切，它迫使灵魂进行哲思，并为来世担心。财富的欲望、获利的爱慕和肉体的享乐不再控制它。这些事如云消散，留下清晰的判断力，悲伤进入，软化了刚硬的心。因为没有什么比享乐的生活更与哲学相悖；另一方面，也没有什么比苦难更有利于哲学。想想贪婪的人那时会怎样。因为，\BibleQuote{一时的苦难，}经上说，\BibleQuote{使人忘记许多享乐。} \BibleRef{德 2:27} 那时他的状态如何，当他想到那些被他抢劫、伤害、欺骗的人，当他看到别人收获他贪婪的果实，而自己要去受罚？因为不可能，真的不可能，当他生病时他不反思这些事。因为灵魂本身常常因痛苦和恐惧而分心。这是何等的苦涩，告诉我！而且每种疾病都必须忍受这些。当他看到别人受惩罚或处死时，他还有什么苦不会受呢？\par

这些事情在此等待着他。至于他日后必须经受的，无法言说有什么惩罚、什么报应、什么折磨、什么苦刑在那里为他存留。我们宣告这些事。\Quote{有耳听的，就应当听！}\BibleRef{路 8:8}我们不停谈论这些，并非出于自愿，而是不得已。因为我们宁愿没有义务提及这些事情。但既然必须如此，我们至少要用一点药物，使你们脱离疾病，恢复健康。但你们仍处于这种病态中时，放弃治疗便是卑劣懦弱的表现，更不必说残忍与不人道了。因为如果当医生对我们的身体绝望时，我们恳求他们不要忽视我们，直到我们最后一息都要尽其所能，那么我们岂不更当劝勉自己？因为也许当我们到了地狱的门口，罪恶的门槛时，仍有可能康复，重获力量，把握永生！有多少人听了十次仍无动于衷，后来却因一次听闻而皈依！更确切地说，并非一次听闻；因为他们虽然在那十次讲道中似乎麻木，却仍有所得，随后一下子结出丰硕的果实。如同一棵树，可能受了十次砍伐而不倒，然后被一击彻底倒下；但这并非那一击之功，而是那十次砍伐使最后那击成功。这一点，看见树根的人知道，而只看见树干外表的人却不知道。这里的情形也是如此。同样，医生用了许多药物，却未见效；后来有人进来，完全治愈了。但这不仅是他一个人的功劳，也是那些已经减轻病情的人之功。因此，即使我们现在没有结出听道的果实，将来也必会结出。因为我相信我们必会结出果实。因为，这般热切的渴望、这般爱听的心，不可能不产生效果。愿天主不许！但愿我们所有人都配得上基督的训诲，获得永恒的福乐，等等。\par

\SourceRecord{Translated by Philip Schaff.；Nicene and Post-Nicene Fathers, First Series；Vol. 13.；Edited by Philip Schaff.；Buffalo, NY: Christian Literature Publishing Co.,；1889.；<http://www.newadvent.org/fathers/230705.htm>.}

% structure: p0001 source=h1 target=section reason=page-title

\section{论弟茂德后书第六篇讲道}

\BibleRef{弟后 2:20-21}\par

\BibleQuote{但在大宅里，不仅有金器银器，也有木器瓦器；有作贵重的，有作卑贱的。所以人若自洁，离开这些，他必成为贵重的器皿，成圣，合主使用，预备行各样的善事。}\par

甚至现在仍有许多人对以下事实感到困惑：为何恶人被容许存留，尚未被消灭。无疑，对此可以提出各种理由，例如：他们可能会悔改，或者他们的惩罚可成为大众的鉴戒。但\Person{保禄}在这里提到一个类似的情况。因为他说道：\par

\BibleQuote{在大宅里，不仅有金器银器，也有木器瓦器。}借此表明，正如在大宅里很可能有器皿的巨大差异，在这里，在整个世界里也是如此——因为他说的不是教会，而是整个世界的范围。因为请想想，我求你们，他并非指教会；在那里他不会有任何木器或瓦器，而是全都是金器或银器，那里有基督的身体，那里是\BibleQuote{无玷无瑕的贞洁童女}。\BibleRef{弗 5:27}他想要说的是：不要让有败坏和邪恶之人这件事困扰你。因为在大宅里就有这样的器皿。但然后呢？他们并没有得到同样的尊荣。有的是为尊贵，有的是为卑贱。\Quote{不然，}有人说，\Quote{在屋里它们或许有些用处，但在世界里毫无用处。}虽然天主不用它们做如此尊贵的服务，却利用它们做其他用途。例如，虚荣的人建造很多，贪婪的人也如此，商人、工匠、官吏；世界上有某些工作适合他们。但金器不是这样的性质；它用于君王的餐桌。然而他并不是说邪恶是必要的（怎么可能呢？），而是说邪恶之人也有他们的工作。因为如果全都是金器或银器，就不需要那些卑贱的种类了。例如，如果所有人都刚强，就不需要房屋了；如果所有人都没有奢侈，就不需要美味了；如果所有人都只关心必需品，就不需要辉煌的建筑了。\par

\BibleQuote{所以人若自洁，离开这些，他必成为贵重的器皿，成圣。}你看见吗？成为金器或瓦器不是出于本性，也不是出于物质的必然性，而是出于我们自己的选择？否则，瓦器不能成为金器，金器也不能降格为另一者的卑贱。但在此例中，有巨大的变化和状态的改变。\Person{保禄}曾是瓦器，却成了金器；\Person{犹达斯}曾是金器，却成了瓦器。所以，瓦器是由于不洁而成为的；行淫的和贪婪的人成为瓦器。\Quote{但他为何在别处说：\InnerQuote{我们有这宝贝在瓦器里}？如此他并非藐视而是尊崇瓦器，视之为宝贝的容器。}那里他显示的是材质本身，而非材料的形式。因为他要说的是，我们的身体是瓦器。正如陶器不过是烘干的泥土，我们的身体也不过是泥土因灵魂的热力而凝结；因为它是瓦器的，是明显的。因为正如这样的器皿常因摔倒而破碎粉碎，我们的身体也因死亡而倒下并分解。因为我们的骨头与陶片有何不同？它们又硬又干。我们的肉体与泥土有何不同？由水组成，正如泥土一样。但是，正如我所说，他为何不轻蔑地说它？因为那里他是在谈论它的本性，这里是在谈论我们的选择。\Quote{人若自洁，}他说，\InnerQuote{不是仅仅洁净，而是洁净出去，}即完全洁净自己，\Quote{他必成为贵重的器皿，成圣，合主使用。}因此，其他器皿对任何善用都是无用的，虽然对他们还作某种使用。\Quote{预备行各样的善事。}即使他没有行，他也适合并有能力行。因此我们应当预备好一切，甚至死亡、殉道、守贞的生活，或所有这些。\par

第22节。\Quote{也要逃避少年人的私欲。}\BibleRef{弟后 2:22}\par

不仅是行淫的私欲，而且一切无节制的欲望都是少年人的私欲。让老年人知道，他们不应该行少年人的事。如果有人蛮横无礼、贪权、贪财、贪肉体享乐，这是少年人的私欲，是愚昧的。这些事情必定出于尚未坚固的心、尚未深深扎根而是游移不定的心思。那么他建议什么，好使人不被这些事物所俘虏？\Quote{要逃避少年人的私欲，}\par

\Quote{要追求正义、信德、爱德、平安，与那些以纯洁的心呼求主的人一起。}\BibleRef{弟后 2:22}\par

他称美德总称为\Quote{正义}：生活的虔诚为\Quote{信德、温良、爱德}。\par

\Quote{那些以纯洁的心呼求主的人}是什么意思？好像是说：不要因那些仅仅呼求主的人而喜乐；而要因那些真诚无伪地呼求主的人、毫无诡诈、在平安中亲近祂、不争辩的人。要与这些人结交。至于其他人，不要轻易与他们交往，但尽可能与众人和睦。\par

第23节。\Quote{但要避免愚拙无知的辩论，因为知道它们能引起争竞。}\BibleRef{弟后 2:23}\par

你看见吗？他处处引领他离开辩论；并不是因为他不能驳倒它们，因为他完全能够。如果他不能，他应该说：要殷勤，使你能驳倒它们；正如他说：\Quote{你要专心读经，因为这样做，你既能救自己，也能救那些听你的人。}\BibleRef{弟前 4:13}\par

第24节。\Quote{主的仆人不应当争辩。}\BibleRef{弟后 2:24}\par

他甚至不应在争论中争辩，因为主的仆人必须远离纷争，既然天主是和平的天主，和平天主的仆人岂能与纷争有干？\par

\Quote{但要对众人温和。}\par

那么他为何说：\Quote{你要以全权谴责}\BibleRef{铎 2:15}，又说：\Quote{不要让人轻视你的青年}\BibleRef{弟前 4:12}，以及\Quote{要严厉地谴责他们}\BibleRef{铎 1:13}？这是因为这与温和并不矛盾。因为严厉的谴责，若以温和的态度发出，往往更能深入创伤：因为以温和触动人心，有时比以威势慑服更为有效。\par

\Quote{善于教导}；即那些愿意受教的人。因为他说：\Quote{异端的人，在第一次和第二次劝诫以后，就应该弃绝。}\BibleRef{铎 3:10}\Quote{忍耐。}他加上这一点很恰当，因为这是教师尤其应当具备的品质。没有忍耐，一切都属徒然。如果渔夫不失望，虽然他们常常整日撒网一无所获，我们更不应失望。请看结果：借着不断的教导，言语的犁常常深入灵魂的深处，拔除那困扰它的邪恶情欲。因为常常听的人终会被触动。人不可能持续听而不产生任何效果。因此，有时他几乎要信服了，却因我们的厌烦而丧失。这正像一个不善的农夫，第一年为他所栽种的葡萄树松土，第二年、第三年都想收获果实，却一无所获；三年后他失望了，而在第四年，当他即将获得劳苦的报酬时，却放弃了葡萄树。说了\Quote{忍耐}之后，他还不满足，继续说道：\par

第25节。\Quote{用温柔教导那些反对自己的人。}\BibleRef{弟后 2:25}\par

因为教导的人尤其应当注意以温和进行。因为一个愿意学习的灵魂，不能从严厉和争论中获得任何有益的教训。当它想要应用时，若被如此搅扰，便什么也学不到。凡愿意获得有用知识的人，首先应对自己的老师怀有好感；若不能先达到这一点，便无法完成任何必要或有益的事。没有人会对粗暴专横的人怀有好感。那么他为何说：\Quote{异端的人，在第一次和第二次劝诫以后，就应该弃绝}？他是指那不可救药的人，是他知道已病入膏肓的人。\par

\Quote{或许天主会赐给他们悔改，以认识真理。}\par

第26节。\Quote{并使它们从魔鬼的网罗中恢复自己。}\BibleRef{弟后 2:26}\par

他所说无非是：或许会有改过迁善。或许！因为这是不确定的。所以我们只应离开那些人，即我们能明确表明、且深信无论做什么都不会悔改的人。\Quote{用温柔，}他说。你看，我们应当以这种态度对待那些愿意学习的人，并与他们不断交谈，直到我们达到论证的地步。\par

\Quote{他们已被他俘虏，顺从他的旨意。}这话说得对：\Quote{他们已被俘虏}，因为在此期间他们漂浮在错误中。请注意他如何教导人要谦卑。他没有说\Quote{或许你能}，而是说\Quote{或许天主会赐给他们恢复}；所以，若成就了什么，一切都是出于主。你栽种，你浇水，但他播种并使之结果。所以我们不要如此受影响，仿佛我们亲自成就了说服，即使我们说服了任何人。\Quote{被他俘虏，}他说，\Quote{顺从他的旨意。}没有人会说这是指教义，而是指生活。因为\Quote{他的旨意}是要我们生活正直。但有些人因他们的生活而陷于魔鬼的网罗，所以我们也不应对这些人感到厌倦。\par

\Quote{或许，}他说，\Quote{那些被俘的人能够恢复，顺从他的旨意。}现在\Quote{或许}表明极大的宽忍。因为不行天主的旨意，就是魔鬼的网罗。\par

因为像一只麻雀，虽未被完全关闭，只被缠住脚，仍受设网者的控制；所以，即使我们未在信仰和生活上完全颠覆，只在生活上，我们仍受魔鬼的控制。因为\Quote{不是凡向我说\InnerQuote{主啊，主啊}的人，就能进入天国}；又说：\Quote{我不认识你们，你们这些作恶的人，离开我吧。}\BibleRef{玛 7:21}你看，当主不认识我们时，信仰毫无益处；他对童女们也说了同样的话：\Quote{我不认识你们。}\BibleRef{玛 25:12}那么童贞或许多劳苦有何益处，当主不认识我们时？在许多地方，我们发现人并非因信仰受责，而是因恶行受罚；在别处，并非因恶行受责，而是因教义败坏而丧亡。因为这些事是相连的。你看，当我们不行天主的旨意时，我们就在魔鬼的网罗之下。而且常常不仅因坏生活，而且因一个缺陷，我们就进入地狱，那里没有善行来平衡它，因为那些童女没有被指控行淫或奸淫，没有嫉妒或恶意，没有醉酒，也没有信仰不纯，而是因为油不够，即缺少施舍，因为那就是所指的油。而那些被咒诅的人们，在\Quote{可咒骂的，离开我，到永火里去}这些话语中，并没有被指控任何此类罪行，而是因为她们没有喂养基督。\par

道德训诲。你看见吗？缺乏施舍足以将人投入地狱之火？因为那不行施舍的人将在何处得救？你每日禁食吗？那些童贞女也如此行，但于她们无益。你祈祷吗？那又怎样？没有施舍的祈祷是不结果实的，没有施舍，一切都是不洁和无益的。德行中更好的部分被摧毁了。经上说：\Quote{那不爱自己弟兄的，}经上说，\Quote{不认识天主。}\BibleRef{若一 4:8}而当你甚至不把那些卑微无价值的东西分给他时，你怎能爱他呢？告诉我，因此，你守贞洁吗？出于什么原因？出于对惩罚的恐惧吗？决不是。你守贞洁是由于本性禀赋，因为如果你出于对惩罚的恐惧而守贞洁，并且违逆本性服从如此严厉的规矩，那么你更应该行施舍。因为克制对财富和肉体快乐的欲望并不是一回事。后者更难约束。为什么？因为这种快乐是本性的，其欲望是与生俱来的，并在身体中自然生长。财富却不然。在此事上，我们能够肖似天主：显示慈悲和怜悯。因此，当我们没有这种品质时，我们就毫无善。祂没有说：\Quote{你们若禁食，就会像你们的父，}也没有说\Quote{你们若守童贞，}或\Quote{你们若祈祷，}也没有说\Quote{你们就会像你们的父，}因为这些事中没有一件能应用于天主，也不是祂的行动。但祂说什么？\Quote{你们应当慈悲，就像你们的天父是慈悲的一样。}\BibleRef{路 6:36}这是天主的作为。因此，如果你没有这个，你有什么呢？祂说：\Quote{我喜欢慈悲，而不是祭献。}\BibleRef{欧 6:6}天主创造了天、地和海洋。这些是伟大的工程，配得上祂的智慧！但祂藉着慈悲和对人类的爱，如此强有力地吸引了人性归向祂，远胜过藉着其它一切。因为那本是能力、智慧和良善的工程。但祂成为仆人，更是如此。我们不是为此更钦佩祂吗？我们不是为此更惊叹于祂吗？没有什么像慈悲那样吸引天主亲近我们。先知们从始至终都论述这个主题。但我所说的不是伴随着贪婪的慈悲。那不是慈悲。因为产生油的不是荆棘之根，而是橄榄之根；同样，产生慈悲的不是贪婪、不义或掠夺之根。不要诋毁施舍。不要使施舍被众人所指责。如果你为此而行抢劫，以便行施舍，那么你的施舍比什么都邪恶。因为它出自掠夺，那便不是施舍，而是不人道，是残忍，是对天主的侮辱。如果加因因奉献自己较次的礼物而得罪，那么奉献他人财物的人岂不更得罪？奉献无非是祭献，是净化，而非玷污。你尚且不敢用不洁的手祈祷，却献上抢劫来的污秽，认为你没有做错吗？你不让你的手充满污秽和脏物，而是先洗净它们，然后奉献。然而那种污秽并不归咎于你，而另一种则应受责备和谴责。因此，不要让我们考虑如何用洁净的手献上祈祷和祭品，而要考虑所献之物是否纯洁。如果有人洗净器皿，却装满不洁的礼物，岂不是可笑的嘲讽？让手洁净吧；它们将会洁净，如果我们不仅用水洗，而且首先用正义洗。这是手的净化者。但如果双手充满不义，即使洗上一千次，也无济于事。祂说：\Quote{你们要洗濯，自洁，}\BibleRef{依 1:16}但祂又说\Quote{去到浴池、湖泊、河流吗？}不。而是什么？\Quote{从你们的心灵除去你们的恶行。}这才是洁净。这才是从玷污中得洁净。这才是真正的纯洁。另一种用处甚微；而这一种赐给我们对天主的信心。前者可为奸夫、盗贼、凶手、卑劣、放荡、柔靡之人所获得，尤其是后者。因为他们常常注意身体的洁净，用香料熏香，清洗他们的坟墓。因为他们的身体不过是坟墓，因为灵魂已死于其中。因此，这种洁净可能为他们所有，但内在的洁净却没有。\par

洗涤身体并无大事。那是犹太式的净化，无意义且无益，因为内心缺乏纯洁。假设一个人患有腐烂的疮口或消耗性的溃疡；无论他怎样清洗身体，都无济于事。如果身体的腐烂不能通过清洁和掩饰外表而受益；那么当灵魂被腐朽感染时，身体的纯洁又有何用？毫无用处！我们的祈祷应当是纯洁的，如果它们从腐朽的灵魂发出，就不能纯洁，而没有什么比贪婪和掠夺更腐蚀灵魂了。但有些人白天犯下无数罪过，晚上洗澡，然后进入教堂，满怀信心地举起双手，仿佛藉着沐浴的洗涤，他们已经摆脱了所有的罪责。如果真是这样，每天洗澡将大有裨益！我本人也不会停止频繁去浴池，如果它能使我们纯洁，洗净我们的罪！但这些事是琐碎可笑的，是儿童的玩具。天主所厌恶的不是身体的污秽，而是灵魂的不洁。因为祂说：\Quote{纯洁的人是有福的}——祂是说身体吗？不——\Quote{心里洁净的人是有福的，因为他们要看见天主。}\BibleRef{玛 5:8}先知又说：\Quote{天主，求祢在我内造一颗纯洁的心。}\BibleRef{咏 50:10}又说：\Quote{洗净我内心脱离邪恶。}\BibleRef{耶 4:14}\par

养成行善的习惯大有益处。请看这些洗濯是何等琐碎而无益。但当灵魂为某种习惯所占据时，便不会离弃它，也不敢在未履行这些仪式之前就前来祈祷。例如，我们养成了洗濯和祈祷的习惯，不洗濯就认为不宜祈祷。我们不愿以未洗的手祈祷，仿佛这样会得罪天主、违背我们的良心。如今，如果这琐细的习俗对我们有如此大的力量，且日日遵守；倘若我们养成了施舍的习惯，并决意恒常遵守，以至永不空手进入祈祷之所，那目的便达到了。习惯的力量是巨大的，无论善恶，一旦习惯推动我们，便没有多少困难。许多人养成了不断画十字的习惯，无需他人提醒，常常当心思游荡于别处时，手便不由自主地因习惯，如同有一位活生生的导师，做出这个记号。有些人养成了绝不发誓的习惯，因此无论情愿与否，他们从不这样做。那么，让我们养成这样的施舍习惯吧。\par

发现如此的补救方法，值得我们付出何等辛劳？且说，假如没有施舍的缓解，而我们仍因无数的罪恶使自己应受天主的惩罚，岂不是要悲哀痛哭？岂不是要说：愿能藉财富洗去我们的罪孽，我们愿舍弃一切！愿能藉财富消除天主的义怒，我们必不惜家产？因为我们在患病时，临终之际也会说：\Quote{若能用钱买免死亡，某某必献出他所有的财产}；更何况在此事上。请看天主对人的爱何等伟大！祂赐我们能力买免的不是暂时的死亡，而是永久的死亡。祂说：不要买这短暂的生命，而要买那永恒的生命。我卖给你的是那一个，不是这一个：我不戏弄你。即使你赢得了今世生命，也不过是虚无。我知道我所赐予你的价值。买卖世间货物的人并不如此行事。他们能欺骗谁时，就付出少许获得巨利。天主并非如此。祂用较少的代价赐下多得多的一切。\par

告诉我，假如你去找一个商人，他摆出两块石头在你面前，一块价值微小，另一块非常宝贵，必定能带来大量财富；如果他允许你以廉价石头的价格拿走那块更贵重的，你会抱怨他吗？不会！你反倒会赞赏他的慷慨。同样，现在有两种生命摆在我们面前，一种是暂时的，另一种是永恒的。天主将这些卖给我们，但祂愿意卖给我们后者而非前者。我们为何像愚蠢的孩子一样抱怨，说我们得到了更珍贵的？难道可以用金钱购买生命吗？是的，当我们所施舍的是自己的财产，而非他人的；当我们不搞欺骗时。但你说，从那以后财产就是我的了。它们在被抢夺后已不再属于你。它们仍是邻人的，即使你千百次地成为其主人。因为若你接收一笔存款，即使在存款人旅行的短暂期间，它也不是你自己的，虽然可能存放在你那里。因此，如果那经存款人同意和感谢才被我们接收的，即使在我们保留的短期内也不是我们自己的，那么那违背主人意愿被我们掠夺的，就更不是我们的了。无论你扣留多久，他才是主人。但德行真正是我们自己的；至于金钱，连我们自己的严格来说也不是我们的，更不用说别人的了。今天属我们，明天就属于另一个人。德行之物才是我们自己的产业。这不会像其他东西那样遭受损失，而是完全被所有拥有它的人所拥有。因此，让我们获取它，让我们轻视财富，好使我们能获得那些真实的财富，愿天主使我们堪当分享这些财富，藉着恩宠与慈爱等等。\par

\SourceRecord{Translated by Philip Schaff.；Nicene and Post-Nicene Fathers, First Series；Vol. 13.；Edited by Philip Schaff.；Buffalo, NY: Christian Literature Publishing Co.,；1889.；<http://www.newadvent.org/fathers/230706.htm>.}

% structure: p0001 source=h1 target=section reason=page-title

\section{弟茂德后书讲道词 第七篇}

\BibleRef{弟后 3:1}\par

\BibleQuote{你也要知道，在最后的日子里，艰难的时期会来到。因为人们将只爱自己、贪婪、自夸、高傲、亵渎、忤逆父母、忘恩负义、不圣洁、没有亲情、不肯和解、诬告、不能自制、凶暴、不爱好善、背信、鲁莽、自大、爱享乐胜过爱天主；有虔诚的外貌，却否认其能力；你要远离这样的人。因为这类人潜入人家，虏获那些满身罪恶、被各种私欲所牵引的轻浮妇女；她们常常学习，却始终不能认识真理。}\par

他在前书曾说：\BibleQuote{圣神明明地说：在后期有些人会背弃信德}\BibleRef{弟前 4:1}；又在本书信别处预言类似的事将要发生；在此又同样地说：\BibleQuote{你要知道，在最后的日子里，艰难的时期会来到。} 他不仅从未来，也从过去预言这事：\Quote{如同雅乃斯和杨布勒曾抗拒梅瑟。} 又从推理：\Quote{在大户人家，不仅有金器银器。} 但他为什么这样做？为使\Person{弟茂德}或我们任何人，在遇到恶人时不致惊扰。如果在\Person{梅瑟}时代有这样的人，以后也会有，那么在现今有这样的人也就不足为奇了。\par

\BibleQuote{在最后的日子里，艰难的时期会来到，}他说，就是极其恶劣的时期。时期怎么会是艰难的？他这样说并不是责怪日子或时期，而是责怪那些时期的人。因为我们习惯这样谈论好时期或坏时期，根据发生在其中的事件，这些事件是由人造成的。他立即列出了根源和源头，由此产生这些及其他一切邪恶，即自大。被这激情控制的人甚至不关心自己的利益。因为当一个人忽视邻人的利益，不关心时，他怎么会关心自己的利益？正如关心邻人事务的人会因此妥善安排自己的事务，所以轻视邻人事务的人也会忽视自己的。因为如果我们彼此是肢体，邻人的福利不仅是他自己的事，也是整个身体的事；邻人的伤害不仅限于他，也会用痛苦搅扰其他一切。如果我们是一座建筑，任何部分薄弱都会影响整体，而坚固的部分则给其余部分带来力量和支持。同样在教会里，如果你轻视你的邻人，你就伤害了自己。怎么伤害？因为你自己的一个肢体遭受了不小的伤害。如果那不分施自己财物的人会下地狱，那么看到邻人遭受更严重的痛苦却不肯伸手相助的人，就更应受谴责，因为在这种情况下损失更严重。\par

\Quote{因为人们将只爱自己。}爱自己的人可以说并不爱自己，而是爱自己弟兄的人才是真正爱自己。从自爱产生贪婪。\Quote{贪婪。}从贪婪产生自夸，从自夸产生骄傲，从骄傲产生亵渎，从亵渎产生悖逆和不服从。因为高抬自己的人对人容易，对天主也容易。这样，罪恶就产生了。它们常常从下层上升。对人虔敬的人，对天主更虔敬。对同伴温和的人，对主人更温和。轻视同伴的人，最终会轻视天主本身。伦理教训。那么我们不要彼此轻视，因为这是一种邪恶的训练，教我们轻视天主。确实，彼此轻视实际上就是轻视天主，因为天主命令我们彼此尊重。这可以用一个例子来证明。\Person{加音}轻视了他的弟弟，随即轻视了天主。怎么轻视？注意他傲慢地回答天主：\BibleQuote{我岂是看守我弟弟的吗？}\BibleRef{创 4:9}同样，\Person{厄撒乌}轻视了他的弟弟，他也轻视了天主。因此天主说：\BibleQuote{我爱了雅各伯，却恨了厄撒乌。}\BibleRef{罗 9:13}因此\Person{保禄}说：\BibleQuote{免得有淫乱和亵圣的人，如同厄撒乌一样。}\BibleRef{希 12:16}\Person{若瑟}的兄弟们轻视了他，他们也轻视了天主。以色列人轻视了\Person{梅瑟}，他们也轻视了天主。\Person{厄里}的儿子们轻视了人民，他们也轻视了天主。你也从反面看到吗？\Person{亚巴郎}对侄儿温柔，他顺从天主，如同他对儿子依撒格的行为以及所有其他德行所显明的。同样，\Person{亚伯尔}对弟弟温和，他也对天主虔敬。因此，我们不要彼此轻视，免得也学会轻视天主。让我们彼此尊重，以便也学会尊重天主。对人傲慢的人，也会对天主傲慢。但当贪婪、自私和傲慢聚在一起时，离完全的毁灭还缺什么？一切都腐化了，一股罪恶的洪流冲入。\Quote{忘恩负义，}他说。因为贪婪的人怎能感恩？他对谁会感恩？对谁都不会。他认为所有人都是敌人，想要所有人的财物。即使你把全部财产花在他身上，他也不会感恩。他生气你不再有更多，以便你可以给他。即使你让他成为全世界的主人，他仍然忘恩负义，认为什么也没得到。这欲望无法满足。它是疾病的渴望；疾病渴望的本性就是这样。\par

患热病的人永远不会满足，不断渴望饮水，却永远不能解渴，而是连续忍受干渴；同样，疯狂追求财富的人，他的欲望永远无法满足；无论给他多少，他仍不满足，因此永远不懂感恩。他不会感激那没有给他所愿之多的人，而这没有人能做到。因为他的愿望没有止境，所以他不会感恩。因此，没有人像贪婪的人那样忘恩负义，没有人像爱财的人那样麻木不仁。他是全世界的敌人。他愤怒世上还有他人存在。他宁愿世界变成一片荒原，好让他独占所有人的财产。他生出许多疯狂的幻想。\Quote{但愿发生地震，}他说，\Quote{全城被吞没，唯独我留下，好占有所有人的财富！但愿瘟疫来临，毁灭一切，只留下黄金！但愿海水泛滥或海啸发生！}这就是他的幻想。他从不祈求好事，只祈求地震、雷击、战争、瘟疫之类。那么，你告诉我，你这可怜的人，比任何奴隶更卑贱，如果万物都是黄金，你岂不是被你的黄金毁灭，饿死吗？如果世界被地震吞没，你也会因你的致命欲望而丧亡。因为如果除了你之外没有其他人，生活必需品就会断绝。假设地上其他居民都瞬间消灭，他们的金银自动归你所有（因为这些人的幻想是荒谬和不可能的）。但如果他们的金银、丝绸和金线衣服都到你手里，对你有什么益处呢？死亡只会更确定地临到你，因为没有人为你准备面包或耕种土地；野兽会四处游荡，魔鬼会因恐惧而折磨你的灵魂。现在确实有许多魔鬼占据它，但那时它们会引你绝望，使你立刻陷入毁灭。但你说：\Quote{我希望有人耕种土地，有人准备食物。}那么他们会消耗一些。\Quote{但我希望他们什么也不消耗。}这欲望多么贪得无厌！还有什么比这更可笑？你看到事情不可能吗？他希望有许多人服侍他，却又吝啬他们的一份食物，因为这会减少他的财产！那怎么办？你希望人是石头吗？这全是胡闹；波浪、风暴、巨浪、剧烈的动荡和暴风雨淹没灵魂。它永远饥饿，永远干渴。我们难道不同情他、为他哀伤吗？在身体疾病中，这被认为是最痛苦的之一，医生称之为\OriginalTerm{善饥症}{bulimy}，即人虽然充满，却总是饥饿。灵魂中同样的紊乱不是更可悲吗？因为贪婪是灵魂的病态饥饿，总是填塞，永不满足，仍然渴求。如果必须喝藜芦或忍受任何千倍更糟的事，难道我们不值得欣然承担，以便从这激情中解脱出来吗？任何丰厚的财富都不能填满贪婪的肚腹。难道我们不羞愧吗？这些人竟如此被爱财之心冲昏头脑，而我们却不显出同等热诚去爱天主，不按金银的尊荣去光荣祂？为了金钱，人会不眠不休、跋涉旅行、持续冒险、忍受仇恨敌对，总之付出一切。但为了天主，我们连一句话都不敢说，也不愿招致仇视；如果我们需要帮助受迫害的人，我们就抛弃受害者，回避权势者的仇恨和危险。虽然天主给了我们力量去援助他，却因我们不愿招人仇恨和不满而任其灭亡。许多人还为此辩护说：\Quote{宁可无缘无故被人爱，也不可无缘无故被人恨。}但这是无缘无故遭恨吗？或者，这样的仇恨岂不更好？因为为天主而受恨，比为天主而受爱更好：当我们为天主而被爱时，我们欠这光荣的债；但当我们为祂而被恨时，祂欠我们赏报。爱财者不知爱财的限度，无论多么大；但我们，即使做了很少，就以为已履行了一切。我们爱天主不及他们爱黄金，差得远。他们对黄金的失控的狂热是对他们的严厉控诉。而我们没有为天主如此疯狂，这是对我们的定罪；我们没有将如此之多的爱献给万有的主，如同他们献给尘土一般，因为矿中挖出的黄金并不更好。\par

让我们看看他们的疯狂，并为自己感到羞愧。因为我们虽然不被爱黄金的热情所燃烧，却不热切地向天主祈祷？因为在他们的情形中，人们轻视妻子、儿女、财产和自身的安全，而且当他们不确定是否能增加财产时也是如此。因为常常，在他们希望的正中间，他们同时失去了生命和辛劳。但我们，虽然知道如果我们像应当爱祂那样爱祂，就会得到我们所渴望的，却仍不爱祂，而且在爱近人和爱天主方面都全然冷淡；对天主的爱冷淡，因为对近人的爱冷淡。因为确实，确实不可能，一个没有爱之感觉的人会有任何慷慨或男子气概，因为一切美善的基础正是爱。\Quote{在这上，} 经上说，\BibleQuote{全部法律和先知都系于此。}\BibleRef{玛 22:40} 就像火在森林中常会清除一切，爱的火焰无论在哪里被接，都会烧尽并清除一切有害于神圣收获的东西，使土壤纯淨，适合接收种子。哪里有爱，一切邪恶都被移除。没有爱财——万恶的根源，没有自私；没有夸耀；因为人为何要在朋友面前夸耀？没有什么比爱更使人谦卑。我们为朋友做仆人的事工，并不以为耻；我们甚至感谢有机会服事他们。我们不吝惜财产，常常也不吝惜自身；因为为了所爱的人，有时也会遭遇危险。哪里有真诚的爱，就没有嫉妒，没有毁谤；我们不仅不诽谤朋友，还堵住诽谤者的口。一切都充满温和与柔顺。没有争斗与争执的痕迹。一切都在散发和平。因为\Quote{爱，} 经上说，\BibleQuote{就是法律的满全。}\BibleRef{罗 13:10} 爱没有任何冒犯之处。为什么？因为哪里有了爱，所有贪婪、抢夺、嫉妒、毁谤、傲慢、伪证和谎言的罪都被消除了。人们作伪证是为了抢夺，但没有人会抢夺他所爱的人，反而会把自己的财物给他。因为我们比从他那里接受时更受亏欠。你们这些有朋友的人知道这一点，我指的是真正的朋友，不仅仅名义上的，而是任何按人应当爱的那样去爱的人，任何真正与另一个人相连的人。让那些不知道的人向知道的人学习。\par

现在我要从圣经中给你们举一个友谊的奇妙例子。撒乌耳的儿子约纳堂爱达味，他的灵魂与达味如此紧密相连，以致达味在为他哀悼时说：\BibleQuote{你对我的爱奇妙，超过妇女的爱情。你受伤致死。}\BibleRef{撒下 1:25} 那么，约纳堂嫉妒达味吗？一点也不，尽管他很有理由。为什么？因为从这些事中，他看出王国会从他那里转移给达味，但他却毫无此类感觉。他没有说：\Quote{这就是那剥夺我父国的人}，反而支持达味得到王权；他为朋友不惜牺牲父亲。但不要让任何人认为他是弒父者，因为他没有伤害父亲，而是制止了父亲的不义企图。他与其说伤害，不如说保全了父亲。他不允许父亲进行不义的谋杀。他多次甚至愿意为朋友而死，并且远非指控达味，他反而制止了父亲的指控。他不嫉妒，反而帮助达味得到王国。我为何谈论财富？他甚至为朋友牺牲了自己的生命。为了朋友，他甚至不畏惧父亲，因为父亲怀有不义的企图，但他的良心却远离这一切。这样，正义与友谊结合在一起。\par

约纳堂就是这样。现在让我们看看达味。他没有机会回报恩情，因为他的恩人在达味统治之前就被带走，在达味所服事的人到达他的王国之前就被杀了。那么，让我们看看那位义人如何尽他所能地表达了他的友谊。他说：\BibleQuote{你对我的爱情非常可爱，约纳堂；你受伤致死。}\BibleRef{撒下 1:25} 仅此而已？这固然不是轻微的表示，但他还因着对父亲的恩情的纪念，多次从危险中救出约纳堂的儿子和孙子，并继续扶持和保护他的孩子们，就像对待自己儿子的孩子一样。我希望所有的人都对活人和死人保持这样的友谊。\par

让妇女们听这话（我特别因为她们才提到死去的人），就是那些再婚，并玷污了亡夫之床的妇女，虽然她们曾爱过他。我并非禁止再婚，也不说它是放荡的证明，因为保禄不许我这样说，他堵住我的口，对妇女说：\Quote{她若出嫁，并没有犯罪。}\BibleRef{格前 7:28}然而让我们注意下文：\Quote{但她若能守节，更为有福。}这状态远胜于另一种。为什么？因许多缘故。因为如果完全不嫁比出嫁更好，那么在此情况下就更是如此。你说：\Quote{但有些妇女不能忍受寡居，而陷入许多不幸之中。}是的，因为她们不知道寡居是什么。因为寡居并非不嫁第二次，就像贞洁并非完全不嫁一样。因为正如\Quote{合宜的事}和\Quote{叫你们可以殷勤事奉主，没有分心的事}是一种状态的标记，另一种状态的标记则是孤苦、\Quote{恒切祈祷}、摒弃奢华享乐。因为\Quote{那宴乐度日的，虽活也是死的。}\BibleRef{弟前 5:6}如果守寡之后，你还有同样的排场、同样的炫耀、同样的衣着，如同你丈夫在世时一样，那么你出嫁倒更好。因为不是结合本身可指责，而是随之而来的许多挂虑。但那些没有错的事，你反而不做；而那些并非无关紧要、易受责难的事，你却陷进去了。因此\Quote{有些转去随从了撒殚}，因为她们不能妥当活出寡居的生活。\par

你愿意知道什么是寡妇，和寡妇的尊严吗？请听保禄论及此事的话：\Quote{若她养育儿女，接待旅客，洗圣徒的脚，救济遭难的人，勤行各样善事。}\BibleRef{弟前 5:10}但当你丈夫死后，你仍穿着同样的奢华财富，无怪乎你不能忍受寡居了。因此，请把这财富转移到天上，你就会发现寡居的重担变得可忍受了。但你说：\Quote{如果我有儿女继承父亲的遗产呢？}也要教导他们轻视财富。将你自己的财产转移，只留给他们足够的一份。也教导他们超越财富。但你说：\Quote{除了我的金银之外，我周围还有一群奴隶，被许多事务压垮，我怎能当得起这一切的照顾，而失去丈夫的支持呢？}这不过是借口、托词，这从许多原因可以看出来。因为如果你不贪图财富，也不求增加你现有的资财，你的担子就会轻省。获取财富比照管财富辛苦得多。因此，如果你切断这一件事——积蓄，并用你的财物周济穷人，天主必用祂保护的手托住你。如果你说这话是出于真心想为无父的儿女保存产业，而不是在这托词下被贪婪所控制，那么那鉴察人心的天主知道如何保障他们的财富，就是那命令你养育儿女的那一位。\par

因为凡以施舍建立的屋宇，是不可能遭遇任何灾祸的，确实不能。就算暂时不幸，终究必兴旺。这对于全家胜过枪矛与盾牌。请听魔鬼怎样论约伯说：\Quote{祢不是四面圈上篱笆，围护他和他的家，并他一切所有的吗？}\BibleRef{约 1:10}为何？请听约伯自己说：\Quote{我是瞎子的眼，瘸子的脚。我是孤儿之父。}\BibleRef{约 29:15}正如那不回避他人灾祸的人，即使在自己不幸中也不会受苦，因为他学会了同情；同样，那不肯承担同情之苦的人，必亲身经历一切忧伤。又如身体疾病，如果脚已坏死，手不施同情去清洗伤口、洗净脓水、敷上膏药，手本身也会患同样的病；同样，那在自己未受苦时不肯服侍他人的人，必要承担自己的苦楚。因为那从其他部分蔓延的邪恶必达于此，那时的问题不再是如何服侍他人，而是如何医治和缓解自身的痛苦。此处亦然。那不肯救济他人的人，自己必成受苦者。撒殚说：\Quote{祢圈护了他内外，}我不敢攻击他！但你说：他受了苦难。是的，但这些苦难成了大善的机缘。他的资财加倍，他的赏报增加，他的义德扩展，他的华冠璀璨，他的奖赏荣耀。他属灵和现世的福分都加增了。他失去了儿女，但他没有得回这些儿女，而是得了别的代替他们，而且那些儿女也为复活而保全。若他们被归还，数目会减少，但现在他得了别的代替，祂也要在复活时将他们呈现。这一切临到他，是因为他慷慨施舍。让我们也照样行，好藉着我们的主耶稣基督的恩宠和慈爱，获得同样的赏报。阿们。\par

\SourceRecord{Translated by Philip Schaff.；Nicene and Post-Nicene Fathers, First Series；Vol. 13.；Edited by Philip Schaff.；Buffalo, NY: Christian Literature Publishing Co.,；1889.；<http://www.newadvent.org/fathers/230707.htm>.}

% structure: p0001 source=h1 target=section reason=page-title

\section{弟茂德后书讲道词 第八篇}

\BibleRef{弟后 3:1-4}\par

\BibleRef{弟后 3:1-4} \BibleQuote{你应知道：在末日，危险的时期必将来临。那时人只爱自己，贪恋钱财，自夸，骄傲，亵渎，不孝顺父母，忘恩，不圣洁，无亲情，不和解，诬告，不能自治，凶暴，不爱好善事，背信，冒失，自大，爱享乐胜过爱天主。}\par

若有人现在因异端者的存在而感到困惑，让他记起，从起初就是如此，魔鬼总是用谬误来攻击真理。天主从起初就应许了善，魔鬼也前来应许。天主种植了乐园，魔鬼欺骗说：\BibleQuote{你们将如同天主一样}\BibleRef{创 3:5}。因为他不能在行动上显示什么，便用言语作了更多的应许。这就是欺骗者的性格。此后有加音和亚伯尔，然后有舍特之子与人的女儿，再后有含与耶斐特、亚巴郎与法郎、雅各伯与厄撒乌；如此直到末了，有梅瑟与术士、先知与假先知、宗徒与假宗徒、基督与假基督。当时就是这样，无论之前还是那时。那时有特乌达，有西满，然后有宗徒，又有赫摩革乃和非肋托这一派。简言之，没有什么时候谬误不被竖立起来与真理对立。因此我们不要沮丧。这事将从起初就被预言了。所以他说：\BibleQuote{你应知道：在末日，危险的时期必将来临。那时人只爱自己，贪恋钱财，自夸，骄傲，亵渎，不孝顺父母，忘恩，不圣洁，无亲情。}忘恩的人就是不圣洁的，这很自然，因为一个不对恩人感恩的人，对他人会怎样呢？忘恩的人是背信的人，是无亲情的人。\par

\Quote{诬告，}即诽谤者。因为那些自觉自身没有善行、反而犯许多罪过的人，在诽谤他人品行中找到了安慰。\par

\Quote{不能自治，}指在舌头、食欲及其他一切方面。\par

\Quote{凶暴，}因此他们的不人道和残忍，当一个人贪婪、自私、忘恩、放荡时。\par

\Quote{不爱好善事，背信，冒失。} \Quote{背信，}背叛友谊；\Quote{冒失，}没有稳定；\Quote{自大，}充满傲慢。\Quote{爱享乐胜过爱天主。}\par

\BibleRef{弟后 3:5} \BibleQuote{具有虔敬的外表，却否认了它的能力。}\par

在致罗马人书中，他这样说：\BibleQuote{在法律上有知识和真理的模范}\BibleRef{罗 2:20}，那里他是称赞的口气；但在这里，他将此罪视为超越其他所有缺陷的恶。为什么？因为他用词的意义不同。因为\Quote{肖像}常被用来表示相似，但有时也表示无生命、无价值的东西。他自己在致格林多人书中说：\BibleQuote{男人不应蒙头，因为他是天主的肖像和光荣}\BibleRef{格前 11:7}。但先知说：\BibleQuote{人只像幻影那样行走}\BibleRef{咏 38:9}。圣经有时用狮子代表王权，如：\BibleQuote{他卧如公狮，如母狮，谁能惊起他？}\BibleRef{创 49:9}，有时表示掠夺，如：\BibleQuote{像咆哮怒吼的狮子}\BibleRef{咏 21:13}。我们自己也是如此。因为事物本身是复杂多变的，它们适合被引用于各种意象和例子。就像我们赞美一位美女时，说她像一幅画；而当赞赏一幅画时，又说它栩栩如生、似乎在呼吸。我们并非表达同一事物，而是前者表示相似，后者表示美丽。这里对于\Quote{形式}也是如此，一处经文指模型、代表、教义、虔敬的范本，另一处则指无生命之物、单纯外表、表象和伪善。因此，没有行为的信德，恰被称为没有能力的形式。因为一个美丽强健的身体，若没有力量，就如同画中的人物；同样，正确的信德而不伴有行为。假设某人\Quote{贪恋钱财、背信、冒失，}却相信正确，那有何益处？如果他缺乏基督徒应有的一切品格，不行虔敬的行为，反而以恶行超越外邦人的不敬虔，成为同伙的祸害，使天主受亵渎，教义因他的恶行被毁谤？\par

\BibleQuote{你要远离这样的人，}他说。但如果在\BibleQuote{末期}人们将如此，这怎么可能？当时大概已有这样的人，至少在某种程度上，尽管没那么极端。但事实上，通过他他警告所有人要远离这种人。\par

\BibleRef{弟后 3:6} \BibleQuote{因为这些人潜入人家，诱惑那些被罪压制、被各种私欲牵引的软弱妇女。}\par

\BibleRef{弟后 3:7} \BibleQuote{她们虽然常常学习，却总无法认识真理。}\par

你看见他们运用那古旧欺骗者的诡计，就是魔鬼用来攻击亚当的武器吗？他说：\Quote{进入人的家}。请注意，他以此说法表明他们的厚颜无耻、卑鄙行径和诡诈。\Quote{诱惑愚昧的妇人}，意思是，那容易受骗的人是\Quote{愚昧的妇人}，毫无男子气概：因为受骗是愚昧妇人的事。\Quote{被罪恶重压}。看从哪里产生他们的说服力：来自他们的罪，来自他们自己意识不到任何善！他说\Quote{重压}非常恰当。因为这一词表明他们罪恶的众多，以及他们的混乱无序状态；\Quote{被各种情欲牵引}。他并没有指责本性，因为他责备的不是一般的妇人，而是这样的妇人。为什么说\Quote{各种情欲}？由此暗示他们的种种缺陷：他们的奢侈、他们的无序行为、他们的放荡。他说\Quote{多样的情欲}，即对荣誉、财富、快乐、任性、名誉的欲望；或许还有其他卑劣的欲望隐含其中。\par

\Quote{常常学习，却永不能认识真理。}他这样说并非为了原谅，而是为了严厉地威胁他们；因为他们的理智麻木了，因为他们用情欲和罪恶压垮了自己。\par

第8节。\Quote{现在，雅乃和杨布勒怎样反对梅瑟，这些人也怎样反对真理。}\BibleRef{弟后 3:8}\par

这些是谁？是梅瑟时代的术士。但这些名字怎么没有在其他地方出现？或是出于传统流传，或是保禄藉着启示知道他们。\par

他说：\Quote{心术败坏的人，于信德是可弃绝的。}\par

第9节。\Quote{但他们不能再有进境，因为他们的愚昧必像那两人的愚昧一样，显明在众人面前。}\BibleRef{弟后 3:9}\par

\Quote{他们不能再有进境}；那么，怎么他在别处说\Quote{他们必更趋于不敬}？\BibleRef{弟后 2:16}他那里是指，他们开始创新和欺骗，就不会在错误中停顿，而总会发明新的骗术和败坏的道理，因为错误从不停滞。但这里他说，他们不能欺骗，也不能引诱人跟他们走，因为无论起初他们怎样貌似能欺骗人，不久便容易被人识破。因为他提到这点，可从下文看出：\Quote{他们的愚昧必显明给众人。}从何而显？处处——\Quote{像那两人的愚昧一样。}因为如果错误起初兴盛，它们不会持续到底，因为非本性真实而只是外表好看的事物就是这样；它们兴盛一时，然后被识破，归于乌有。但我们的道理却不是这样，你便是见证，因为在我们的道理中没有欺骗，因为谁会为欺骗而选择死亡呢？\par

\Quote{至于你，你却追随了我的教训。}\BibleRef{弟后 3:10}因此要坚强；因为你不仅在场，而且紧紧跟随。这里他似乎暗示时期已经很长，因为他说\Quote{你追随了我的教训}；这是指他的言论。\Quote{我的品行}；这是指他的行为。\Quote{我的志向}；这是指他的热忱和灵魂的坚定。他说，我并非只说这些而不实行；也不是只在言语上作哲人。\Quote{信德、坚忍}。他意思是，这些事没有一件困扰我。\Quote{爱德}，这是那些人所没有的；\Quote{忍耐}，他们也没有此项。他意思是，我对异端者显出了极大的坚忍；\Quote{忍耐}，是指在迫害中的忍耐。\par

\Quote{迫害、苦难。}\BibleRef{弟后 3:11}\par

有两件事困扰教师：异端者的众多，以及人缺乏忍受苦难的毅力。然而关于这些，他已经说了很多：这样的情形一直有，也将永远有，没有时代能免于它们，并且它们不能伤害我们，在世界上有金器银器。你看他如何继续谈论他的苦难，\Quote{就是我在安提约基雅、依科尼雍、吕斯特辣所受的。}\par

为什么他从众多事例中选择了这几件？因为其余的事弟茂德已经知道，而这些或许是新近发生的事，他没有提及以前的事，因为他不是逐一举述，也不是出于野心或虚荣，而是为了安慰他的门徒，并非炫耀。这里他指的是丕息狄雅的安提约基雅和吕斯特辣，即弟茂德本人的家乡。\Quote{我所受的迫害。}这里有双重的安慰：我显示了慷慨的热忱，并且我没有被遗弃。不能说天主舍弃了我，反而祂使我的荣冠更加灿烂。\par

\Quote{我所受的迫害，但主从这一切中拯救了我。}\par

\Quote{凡愿意在基督耶稣内虔诚生活的人，都必受迫害。}\BibleRef{弟后 3:12}\par

但是为什么他只说我自己的事？他说：每一个愿意虔诚生活的人都会受迫害。这里他把苦难和忧苦称为\Quote{迫害}，因为一个走德行之路的人不可能不遭受忧伤、苦难和诱惑。因为走窄狭小路、又听见\Quote{在世界上你们要受苦难}的人，怎能幸免呢？\BibleRef{若 16:33}如果约伯在他的时代说\Quote{人在世上，是一场考验}，\BibleRef{约 7:1}何况在那时呢？\par

\Quote{可是邪恶的人和欺诈者必越来越坏，他们欺骗人也被人欺骗。}\BibleRef{弟后 3:13}\par

他说：不要让这些事困扰你们，如果他们在顺境，而你们在试炼中。情况本就如此。从我的例子你可以学到，人在与恶人争战时，不可能不遭遇苦难。一个人不能一边作战一边享乐，不能一边摔跤一边宴饮。因此，那些争斗的人不应寻求安逸或欢乐的生活。再者，目前的状态是竞赛、争战、患难、困苦和考验，正是冲突的战场。安息的时候尚未到来，现在是劳苦和辛劳的时候。没有一个刚脱衣抹油的人会想到安逸。如果你想安逸，为何脱衣、准备战斗？你说：\Quote{难道我不在战斗吗？}什么？当你没有征服你的欲望，也没有抵抗本性的邪恶倾向时？\par

\BibleRef{弟后 3:14} \BibleQuote{但你要持守你所学习且确信的事，因为你知道你是从谁学来的；并且你自幼就认识了圣经，这圣经能使你借着在基督耶稣内的信德获得救恩的智慧。}\par

这是什么意思？正如先知达味劝告说：\BibleQuote{不要嫉妒作恶的人}\BibleRef{咏 36:1}，保禄也劝勉说：\BibleQuote{但你要持守你所学习的事}，不仅是学习，而且是\BibleQuote{确信}，即你所相信的。我所相信的是什么？这就是生命。如果你看到事情与你的信仰相反，不要困扰。同样的事也发生在亚巴郎身上，但他并未受影响。他曾听到\BibleQuote{藉依撒格你的后裔将得名}\BibleRef{创 21:12}，并且他被命令祭献依撒格，但他并未困扰或惊慌。不要因恶人而跌倒。这是圣经从一开始就教导的。\par

那么，如果善人在顺境中，恶人受罚，会怎样呢？一种情况可能发生，另一种则不然。因为恶人可能会受罚，但善人不能总是欢喜。没有人与保禄相等，但他却一生在苦难中度过，昼夜流泪叹息。他说：\BibleQuote{有三年之久，我昼夜不停地流泪劝诫你们每一个人}\BibleRef{宗 20:31}。又说：\BibleQuote{那每日压在我身上的事}\BibleRef{格前 11:28}。他并非今日欢喜，明日忧伤，而是每日不停地忧伤。那么他怎么说：\BibleQuote{恶人会变得更坏}？他没有说他们将得安息，而是说\BibleQuote{他们会变得更坏}。他们的进步是向更坏的方向。他没有说他们会处于顺境。但如果他们受罚，他们受罚是为了使你不至于以为他们的罪未受报应。因为既然我们不被地狱的恐惧所阻止，他就在极度的仁慈中唤醒我们的麻木，唤醒我们。如果没有恶人曾受罚，就没有人会相信天主掌管人间事务。如果所有人都受罚，就没有人会期待将来的复活，因为所有人都在这里得到了应得的报应。为此，他既惩罚，也暂缓惩罚。为此，义人在这里遭受患难，因为他们是寄居者、是旅客，身在异乡。因此，义人忍受这些是为了试炼。因为请听天主对约伯所说的话：\BibleQuote{你以为我警告你是为了别的，而是为了显出你是公义的吗？}\BibleRef{约 40:3}但是罪人遭受任何痛苦，都是在受他们罪的惩罚。所以，无论在患难还是顺境中，让我们都感谢天主。因为二者都有益处。他行事毫无怨恨或敌意，而是出于对我们的关怀和照顾。\par

\BibleQuote{并且你知道你自幼就认识了圣经。}他将圣经称为\Quote{神圣的著作}。你是在这些圣经中长大的，因此借着它们，你的信德应当坚定而不动摇。因为根基扎得很深，并由时间滋养，没有什么能颠覆它。\par

论及圣经，他又加上：\BibleQuote{这圣经能使你获得智慧}，即它们不会使你产生任何愚昧的情感，像多数人所具有的那样。因为正确认识圣经的人，不会因所发生的事而跌倒；他勇敢地承受一切，部分归于信德和天主上智安排的不可理解性，部分则知道其中的理由，并从圣经中找到榜样。因为不去追根究底，不愿知道一切，是知识的一大标志。如果你们允许，我将用例子来说明。让我们假设一条河，或者多条河（我不要求你认可，我只说河流的实际情形），它们的深度并不相同。有些河床浅，有些深得足以淹死不熟悉的人。一处有漩涡，另一处没有。因此，不去尝试所有地方是好的，不愿探测所有深度也是知识不小的证明；而那个敢于试遍河流每一处的人，实际上对河流的特性最无知，往往会因以穿越浅滩时的同样勇气冒险进入深水而丧命。同样，在天主的事上也是如此。那个想知道一切，敢于闯入每件事的人，正是最不知道天主是什么的人。关于河流，大部分是安全的，深水和漩涡很少；但关于天主的事，大部分是隐藏的，无法探究他的作为。那么，你为何要这些深水中淹死自己呢？\par

然而，你当知道，天主安排万事，祂供给一切，我们是\OriginalTerm{自由行动者}{free agents}，有些事祂亲自运作，有些事祂允许发生；祂不愿任何邪恶发生；并非所有事都由祂的旨意成就，有些也由我们的旨意成就；一切恶事仅由我们自己的旨意，一切善事则藉我们自己的意志与祂的推动共同成就；没有一件事是祂所不知的。因此，祂运作万事。你既知道这些，便能分辨什么是善、什么是恶、什么是中性的。如此，德行是善，邪恶是恶；而财富与贫穷、生命与死亡，都是中性的事。你若知道这些，便会因此知道：义人受苦，是为得冠冕；恶人受苦，是为受其罪的惩罚。但并非所有罪人都在今生受罚，免得大众不信复活；也并非所有义人都受苦，免得人以为邪恶而非德行蒙悦纳。这些就是规则和界限。你拿任何事来以此检验，就不会被疑惑困扰。因为正如在计算者当中有六千之数，一切数目都可化约为它，一切都可以在六千的尺度上被除和乘，凡精通算术的人都知道这一点；同样，凡知道这些规则（我将简要复述）的人，永不会跌倒。这些规则是什么？就是：德行是善，邪恶是恶；疾病、贫穷、虐待、诬告等类，都是中性的事；义人今世受苦，或者如果他们有时处于顺境，那是为使德行不致显得可憎；恶人今世享乐，好使他们在将来受罚，或者如果他们有时受罚，那是为使邪恶不致显得蒙悦纳，其行为不致不受惩罚；并非所有人都受罚，免得人不相信复活的时候；甚至善人中有些行了恶事，在此世就清了账；恶人中有些人行了善事，在此世就得了赏报，好使他们的邪恶在将来受罚\BibleRef{玛 6:5}；天主的作为大半是不可理解的，我们与祂之间的差别大得无法言说。如果我们按这些理由来推理，就没有什么能困扰或迷惑我们。如果我们不断聆听圣经，必会找到许多这样的例子。\par

他说：\Quote{这些能使你获得智慧，以致得救。}\par

因为圣经指示我们当做之事与不当做之事。请听这位蒙福者在别处所说：\Quote{你自信是盲人的向导、黑暗中人的光明、愚昧人的训导者、孩童的教师。}\BibleRef{罗 2:19} 你看见法律是黑暗中人的光明；如果那展示文字（文字叫人死）的是光，那么赐予生命的圣神又是什么呢？如果旧约是光，那么新约——它包含如此众多而伟大的启示——又是什么？其差别之大，犹如有人为只知地上事物的人开启天堂，使天上的一切都可见。在那里我们得知关于地狱、天堂和审判的事。我们不要相信不合理的事。它们不过是欺骗。你说：\Quote{什么？}你说道，\Quote{当他们预言的事应验时呢？}这是因为你相信它，如果它确实应验的话。骗子已经俘虏了你。你的生命在他手中，他可随意支配你。如果一个强盗头子将一个投奔他的国王之子置于自己的权力和支配之下，这青年宁愿选择荒漠和无法无天的伙伴，那么这强盗头子能否断言他是生是死？他当然能，不是因为他知道未来，而是因为他是这人生死的支配者，那青年已自投于他的权下。因为凭他自己的喜好，他可以杀死他，也可以饶他一命，既然青年已服从于他，他同样可以随意断言你是富还是贫。世上大部分人都将自己交到了魔鬼手中。\par

再者，一个人已习惯于相信这些骗子，这大大助长了他们的伪装。因为没有人注意他们的失败，却观察他们幸运的猜测。但如果这些人有任何预言的能力，把他们带到我——一个信徒——面前。我说这话并非自夸（因为胜过这些事并非什么大光荣），而且我确实深陷罪中；但在这些事情上，我不愿谦卑；因天主的恩宠，我藐视它们全部。把这个伪称会魔法的人带来；让他，如果他有任何预言能力，告诉我明天会发生什么事。但他不会告诉我。因为我是在君王的权下，他对我的忠诚或服从毫无要求。我远离他的洞穴和地窟。我在君王麾下作战。你说：\Quote{但有人偷了东西，}你说道，\Quote{这个人发现了。}这当然不总是真的，但大部分是荒唐和虚假。因为他们一无所知。如果他们真的知道什么，他们倒应该谈论自己的事：他们的偶像有多少供品被盗，他们多少黄金被熔化。为什么他们没有通知他们的司祭？即使为了钱，当他们的偶像庙宇被烧、许多人同归于尽时，他们也无法提供信息。为什么他们不为自己的安全着想？但若他们预测了什么，那完全是出于偶然。在我们这里有先知，他们从不失败。他们并非有时说真话、有时说假话，而是始终宣告真理；因为这就是预知的特权。\par

因我恳求你们，如果你们至少相信基督，就停止这种疯狂吧；如果你们不信，你们为何暴露自身？你们为何自欺？\BibleQuote{你们心持两意要到几时？}\BibleRef{列上 18:21} 你们为何到他们那里去？为何请教他们？你们一到他们那里，一请教他们，就使自己成为他们的奴隶。因为你们请教，好像你们相信似的。你们说：\Quote{不，}\Quote{我并不是因为相信而请教，而是为了试探他们。} 然而，试探他们是否说真话，这不是相信他们虚假之人的做法，而是仍然怀疑之人的做法。那么你们为何探问将要发生的事？因为即使他们回答说：\Quote{这事将要发生，但你若这样做那样做，就能逃脱}，即使在这种情况下，你也绝不该成为拜偶像者；但你的疯狂还不至于如此大。但如果他们预言将来的事，听他们的人将一无所获，只有无用的悲伤。事情并未发生，他却让自己遭受不安和折磨。\par

如果这事为我们有益，天主就不会吝惜不让我们预先知道。那位将天上的事启示给我们的人，不会嫉妒我们。因为，祂说：\BibleQuote{凡我从圣父所听见的，都已经告诉你们}；并且，\BibleQuote{我不再称你们为仆人，而是朋友。你们是我的朋友。}\BibleRef{若 15:15} 那么，祂为什么没有将这些事启示给我们？因为祂不愿我们关心这些事。为证明祂不嫉妒我们知道这事，这些事曾启示给古人，因为他们还是婴儿，甚至关于一头驴之类的事。但对我们，因为祂不愿我们关心这类事，祂就无意启示它们。但我们会学到什么？是古人从未知道的事，因为古时那些事实在微不足道。但所教导我们的是：我们将复活，我们将是不朽的、不能腐朽的，我们的生命没有终结，万物都要消逝，我们将被提到云中，恶人将受惩罚，以及无数其他事，这一切没有虚谎。知道这些，岂不比听说丢失的驴找到了更好？看，你得了你的驴！看，你找到了它！你得到什么好处？它不久不是还会从别的途径丢失吗？因为即使它不离开你，至少在你死时，你也会失去它。但我所提的这些事情，如果我们紧紧抓住，我们将永远保有。因此，让我们追求这些事。让我们依附这些稳固而持久的善。让我们不要听信占卜者、算命先生和变戏法的人，而要听信那确知万事、拥有万有知识的天主。这样，我们将知道一切适合我们知道的事，并将获得一切美善，藉着恩宠和慈爱，等等。\par

\SourceRecord{Translated by Philip Schaff.；Nicene and Post-Nicene Fathers, First Series；Vol. 13.；Edited by Philip Schaff.；Buffalo, NY: Christian Literature Publishing Co.,；1889.；<http://www.newadvent.org/fathers/230708.htm>.}

% structure: p0001 source=h1 target=section reason=page-title

\section{关于弟茂德后书第九篇讲道}

弟茂德后书 3:16, 17\par

\Quote{凡圣经都是天主所默感的，为教训、为督责、为矫正、为教养在义德上都是有益的：为使天主的人成为完备的，装备齐全，以行一切善事。}［修订本：凡天主所默感的圣经，也是有益的，等等。］\par

他从别处提供了许多劝勉和安慰之后，又添加了更完美的，即从圣经而来的；他如此充分地给予安慰，是合理的，因为他要说一件重大而忧伤的事。因为如果厄里叟，他一直与他的老师在一起直到最后一口气，当他看到老师如同死去般离去时，就悲伤地撕裂了衣服，那么你想这位如此热爱又被爱的门徒，听到他的老师即将去世，并且在他临死时不能享受他的陪伴——这比一切都更令人痛苦——他会遭受什么？因为当我们被剥夺了与逝去者最近的交往时，我们对过去时光的感激之情便会减少。因此，当他先前提供了许多安慰之后，便谈论自己的死亡：而且这不是以普通方式，而是用适合安慰他并使他充满喜乐的话语；这样，就把它视为一种祭献，而不是死亡；一种迁移，事实上正是如此，是迁往更好的状态。\Quote{因为我已被奠祭}\BibleRef{弟后 4:6}，他说。为此他写道：\Quote{凡圣经都是天主所默感的，为教训、为督责、为矫正、为教养在义德上都是有益的。}凡什么圣经？他是指所有那些神圣的著作，就是他正在谈论的。这是关于他所论述的；关于这些他说：\Quote{你自幼便通晓圣经。}因此，凡这样的圣经\Quote{都是天主所默感的}；所以，他说，不要怀疑；并且它\Quote{为教训、为督责、为矫正、为教养在义德上都是有益的：为使天主的人成为完备的，装备齐全，以行一切善事。}\par

\Quote{为教训。}因为从那里我们将知道，我们是否应当学习或不知道任何事情。并且从那里我们可以驳斥虚假的，从那里我们可以被纠正并带回到正轨，可以得到安慰和抚慰，如果有什么缺乏，我们可以得到补足。\par

\Quote{为使天主的人成为完备的。}因为圣经的鼓励就是为了这个，使天主的人藉着它得以完备；因此没有它，他不能完备。他说，你有圣经代替我。如果你要学习什么，你可以从它们学习。如果他这样写信给充满圣神的弟茂德，何况对我们呢！\par

\Quote{装备齐全，以行一切善事；}他不仅是指参与其中，而是\Quote{装备齐全。}\par

第四章 1节。\Quote{我在天主面前，并在将来审判活人死人的基督耶稣面前，恳切嘱咐你。}\BibleRef{弟后 4:1}\par

他是指恶人和义人，或者是指已逝者和仍活着的人；因为许多人将仍活着。在前一封书信中，他引发他的畏惧，说：\Quote{我在使万物生活\BibleRef{弟前 6:13}的天主面前嘱咐你}：但在这里，他把更可怕的事摆在他面前：\Quote{将要审判活人死人的基督耶稣}，即，祂将在\Quote{他的显现和他的国度}中叫他们交账。祂何时审判？当祂在荣耀中显现，并在祂的国度中。他说这话，或是要表明祂不会像祂现在来那样再来，或是\Quote{我呼求他的显现和他的国度作证。}他呼求祂作证，表明他使他想起那显现。然后教导他应当如何传道，他补充说：\par

第二节。\Quote{务要传道：无论得时不得时，总要专心；并要以百般的忍耐和各样的教训责备人、警戒人、劝勉人。}\BibleRef{弟后 4:2}\par

\Quote{得时不得时}是什么意思？就是说，不要有任何限定的时节：让任何时候都成为你的时机，不只在平安、安全以及坐在教堂里的时候。无论你是在危险中、在监狱中、在锁链中，或是要去赴死，就在那时责备。不要停止责备，因为责备在此时最合时宜，就是当你的责备最成功的时候，当事实被证明的时候。\Quote{劝勉人}，他说。像医生那样，先显露伤口，然后切开，再敷上膏药。因为如果你省略了其中任何一项，另一项就变得无用。如果你责备而没有使人信服，你会显得鲁莽，没有人会容忍；但在事情被证明之后，他会顺服责备；之前，他会倔强。如果你使人信服并责备，但很严厉，而不施行劝勉，你所有的劳作都会落空。因为信服本身是难以忍受的，若不混合安慰。正如切开虽然本身有益，但如果没有大量缓和剂来缓解疼痛，病人就不能忍受割和凿，在这件事上也是如此。\par

\Quote{以百般的忍耐和各样的教训。}因为责备的人需要忍耐，免得轻率相信，并且责备需要安慰，好使它被恰当地接受。为什么在\Quote{忍耐}之外他加上\Quote{教训}？\Quote{不是出于愤怒，不是出于仇恨，不是凌辱他，不是好像抓住了敌人。愿这些事情远离你。}而是怎样？如同爱他、同情他，比他自己更因他的悲伤而痛苦，为他的痛苦所融化？\Quote{以百般的忍耐和各样的教训。}这里暗示了不平凡的教导。\par

第三节。\Quote{因为时候将到，那时人不接受健全的道理。}\BibleRef{弟后 4:3}\par

在他们变得顽固之前，你当预先占据他们。为此他说\Quote{得时不得时}；作一切事，以便有愿意的弟子。\par

\Quote{他们必随着自己的私欲，增添许多师傅，}他说。\par

没有什么比这些话更有表现力了。因为通过说\Quote{他们要为自己堆积}，他显示了教师们的混杂众多，也因他们被门徒们选立而显示出来。\Quote{他们要为自己堆积教师}，他说，\Quote{耳朵发痒}。寻求那些说话以满足和取悦听众的人。\par

第4节。\BibleQuote{他们要转耳不听真理，而偏向无稽的传说} \BibleRef{弟后 4:4}。\par

他说这话，不是要将他投入绝望，而是要预备他坚定地承受将要发生的事。正如基督也曾说：\BibleQuote{他们会把你们交出，鞭打你们，并带到会堂前，因我的名} \BibleRef{玛 10:17}。这位蒙福的人别处也说：\Quote{我知道，在我离开后，必有凶暴的豺狼进入你们中间，不顾惜羊群} \EditorialNote{511 xx. 29}。但他说这话，是要他们警醒，好好把握现在的机会。\par

第5节。\BibleQuote{你却要凡事谨慎，忍受磨难} \BibleRef{弟后 4:5}。\par

正因如此，他预言了这些事；正如基督在末时也曾预言将有\Quote{假基督和假先知}；同样，当他即将离去时，也说了这些事。\Quote{你却要凡事谨慎，忍受磨难}；也就是说，要劳苦，趁这瘟疫尚未侵袭他们之前，占据他们的心思；在豺狼进入之前确保羊群的安全，处处忍受艰辛。\par

\Quote{作传福音者的工作，充分证明你的职务}。因此，传福音者的工作就是要忍受磨难，无论是在自己身上，还是来自外人；\Quote{充分证明}，即完成\Quote{你的职务}。看，还有另一个他必须忍受磨难的理由，\par

第6节。\BibleQuote{我现在已准备被奠祭，我离世的时候也近了} \BibleRef{弟后 4:6}。\par

他没有说我的祭献；而是说\Quote{我的被奠祭}，这更强烈。因为并非整个牺牲都献给了天主，而是整个奠祭都献上了。\par

第7节。\BibleQuote{那美好的仗我已打了，赛程我已跑完，信德我已保持了} \BibleRef{弟后 4:7}。\par

当我常常把宗徒拿在手中，思考这段经文时，我不明白为甚么保禄在这里说得如此高超：\Quote{那美好的仗我已打了。}但如今靠天主的恩宠，我似乎找到了答案。那么他为什么这样说呢？他渴望安慰他弟子的沮丧心情，因此吩咐他振作起来，因为他即将得到他的冠冕，已完成了一切工作，获得了光荣的结局。你应该喜乐，他说，而不是悲伤。为什么呢？因为\Quote{那美好的仗我已打了。}如同一个父亲，他的儿子坐在他身边，哀叹自己的孤儿处境，他可能安慰他说：不要哭泣，我的儿子；我们度过了一生，我们到了老年，现在离开你。我们的生命无可指摘，我们光荣地离去，你可以因我们的行为受人敬佩。我们的国王亏欠我们很多。好像他曾说，我们竖起了战利品，我们征服了敌人，这不是夸耀。千万不要这样；而是为了扶起他沮丧的儿子，用赞美鼓励他坚定地承受所发生的事，怀抱美好的希望，不认为这是难以忍受的悲痛事。因为离别实在是令人悲伤的；且听听保禄自己说：\BibleRef{得前 2:17}。如果他因与门徒分离而如此难过，那么你想弟茂德的感受是怎样的？如果在他活着时分别之际他哭泣，以至于保禄说：\BibleRef{弟后 1:4}，那么在他死时岂不是更加难过吗？他写这些事是为了安慰他。事实上，整封信函充满了安慰，是一种遗嘱。\Quote{那美好的仗我已打了，赛程我已跑完，信德我已保持了。}\Quote{一场美好的仗，}他说，所以你要投入其中。但哪里是美好的仗，充满监禁、锁链和死亡？是的，他说，因为这是为基督的缘故而战，且赢得巨大的冠冕。\Quote{美好的仗}！没有任何竞赛比这更高贵。这冠冕没有止境。这不是橄榄叶。没有人类的裁判。没有人类作观众。剧场挤满了天使。在那里人们劳苦多日，忍受艰难，只在那一小时得到冠冕，随即所有快乐消失。但这里完全相反，它永远持续在光明、荣耀和尊贵中。从此我们应当喜乐。因为我要进入安息，我离开赛程。你们曾听说\Quote{离开与基督同处更好。}\par

我已跑完了\Quote{赛程}。因为我们既应较量，也应奔跑；较量，是坚定地忍受磨难；奔跑，不是徒然，而是朝向某个美好的终点。这确实是一场美好的仗，不仅使观众愉悦，也使观众得益；赛跑不是毫无结果。它不只是力量和角力的展示。它将一切带向天堂。这赛跑比太阳更明亮，是的，保禄在地上跑的，比他在天上跑的还亮。他怎样\Quote{跑完了他的赛程}？他走遍了全世界，从加里肋亚和阿剌伯开始，直到地极，以致正如他所说：\BibleQuote{从耶路撒冷直到依里黎孔，我已完全传遍了基督的福音} \BibleRef{罗 15:19}。他像鸟一样飞过大地，或者比鸟更快：因为鸟只飞过，但他拥有圣神的翅膀，穿越了无数障碍、危险、死亡和灾难，因此他甚至比鸟更快。若他只是一只鸟，他可能落下被擒，但被圣神高举，他超越所有罗网，如同有火翅膀的鸟。\par

\Quote{我持守了信德，}他说。有许多东西会夺去它，不只是人的友谊，还有威胁、死亡、以及无数其他的危险：但他屹立抵挡一切。如何做到？藉着清醒与醒寤。这也许足够安慰他的门徒，但他进一步加上赏报。这些是什么呢？\par

第8节。\BibleQuote{从今以后，为我存留了正义的冠冕。}\BibleRef{弟后 4:8}\par

这里他再次称一般的德行为义德。你们不应因我将离去而悲伤，为要戴上那将由基督放在我头上的冠冕。但如果我仍留在此处，你们倒更应悲伤，恐怕我会失败而丧亡。\par

第8节。\BibleQuote{这正义的冠冕，主，正义的审判者，要在那一日赐给我；不但赐给我，也赐给凡喜爱祂显现的人。}\BibleRef{弟后 4:8}\par

这里他提升了他的心神。如果\Quote{给众人}，更给弟茂德。但他没有说\Quote{也给你}，而是说\Quote{给众人}；意思是，如果给众人，更给他。\par

伦理。但有人或许会问，人如何能'爱基督的显现'（\OriginalTerm{显现}{\GreekText{τὴν} \GreekText{ἐπιφανειαν}}）？因祂的来临而喜乐；那因祂来临而喜乐的人，会行相称于其喜乐的事；若有需要，他会舍弃自己的财物，甚至生命，以便获得将来的福乐，堪当以相称的状态、以信心、光明和荣耀，看到那第二次来临。这就是'喜爱祂的显现'。那喜爱祂显现的人，必尽一切努力，在祂普遍来临之前，确保祂特别临到自己。人会说，这怎么可能？请听基督的话，祂说：\BibleQuote{若有人爱我，必遵守我的话，我父也必爱他，并且我们要到他那里去，与他同住。}\BibleRef{若 14:23} 并且想一想，那将普遍显现于众人者，竟应许特别临到我们，是何等的特恩：因为祂说：\BibleQuote{我们要到他那里去，与他同住。} 若有人'喜爱祂的显现'，他必尽一切努力邀请祂到自己内，并持守祂，使那光照耀他。愿没有不配祂来临的事，祂必将很快与我们同住。\par

它被称为祂的\OriginalTerm{显现}{Epiphany}，因为祂将从上显现，从高处照耀。因此我们当\Quote{寻求上面的事}，我们就必很快吸引那些光芒临到我们。那些匍匐于下，埋没于这低下的土地的人，无人能看到那太阳的光。那些以世俗之事污秽自己的人，无人能看见那正义的太阳。祂不照耀任何如此忙碌的人。你当稍作恢复，从那深渊，从那世俗生活的波涛中恢复，如果你想见到那太阳并享受祂的显现。那时你将带着极大的信心看见祂。如今就当做一个哲人。不要让乖戾的灵占据你，免得祂严厉击打你，使你降卑。不要使你的心刚硬，也不要黑暗，免得你在那里失事沉船。不要自欺。因为海底的暗礁造成最致命的船难。不要喂养野兽，我指的是邪恶的情欲，它们比野兽更恶劣。不要信赖永远流逝的事物，以便你能坚定站立。无人能立于水上，但在磐石上所有人都能找到稳固的立足。世俗之事如同水，如同急流，转瞬即逝。他说：\BibleQuote{大水进入我的灵魂。}\BibleRef{咏 68:1} 又说：\BibleQuote{祢把我的脚立在磐石上。}\BibleRef{咏 39:2}\par

有人问：我们如何能为基督受苦？若有人无端诬告你，并非因基督的缘故，但你若忍耐承受，若感恩，若为他祈祷，这一切你都是为基督而做。但若你咒骂他，若你口出怨言，若你试图报复，即使你未能做到，那也不是为基督的缘故；你遭受损失，因你的意图而被剥夺了赏报。因为受苦可以为我们带来益处或害处，这取决于我们。它不取决于苦难的性质，而取决于我们内心的倾向。例如，约伯的苦难很大，但他以感恩承受；他被称义，不是因为他受苦，而是因为在受苦中他感恩地承受。另一人在同样的苦难下（尽管不完全相同，因为无人像约伯那样受苦），但在较轻的苦难下却喊叫、急躁、咒骂全世界、抱怨天主。他被定罪、判罚，不是因为他受苦，而是因为他亵渎；他亵渎，并非出于苦难的必然性，因为如果事件导致的必然是原因，那么约伯也必亵渎；但既然他遭受更严重的苦难却没有做这样的事，可见并非由此原因，而是由于人意志的软弱。因此我们需要灵魂的力量，这样就没有什么会显得沉重；但若我们的灵魂软弱，我们就会对一切感到委屈。\par

根据我们的心境，一切事物都变得可以忍受或不可忍受。让我们坚强我们的决心，我们就能轻易承受一切。那根深扎于地下的树，即使风暴再猛烈，也不会动摇；但如果它只是浅浅地种在地表，一阵微风就足以将它连根拔起。我们也是如此：如果我们因对天主的敬畏而将肉体钉住，就没有什么能动摇我们；但如果我们仅有好的意向，一次小小的打击就会颠覆并摧毁我们。因此，我劝勉大家，要怀着极大的喜乐承受一切，效法那先知，他说：\Quote{我的灵魂依附于祢}；请注意，他说的不是接近，而是\Quote{依附于祢}；他又说：\BibleQuote{我的灵魂渴慕祢。}\BibleRef{咏 61:3}他说的不仅仅是\Quote{渴望}，而是用这样的词语来表达他渴望的热切程度；又说：\BibleQuote{使我的肉身扎根于对祢的敬畏。}\BibleRef{咏 117:120}因为他希望我们如此依附并与他结合，以致永远不与他分离。如果我们这样持守天主，如果我们这样将思念钉在他身上，如果我们因爱他而渴慕，那么我们所渴望的一切都将属于我们，我们也将获得将来的美善，在我们的主耶稣基督内——愿光荣、权能与尊崇归于他，及圣父与圣神，自今以至永远。阿们。\par

\SourceRecord{Translated by Philip Schaff.；Nicene and Post-Nicene Fathers, First Series；Vol. 13.；Edited by Philip Schaff.；Buffalo, NY: Christian Literature Publishing Co.,；1889.；<http://www.newadvent.org/fathers/230709.htm>.}

% structure: p0001 source=h1 target=section reason=page-title

\section{弟茂德后书讲道辞第十篇}

弟茂德后书 4:9-13\BibleRef{弟后 4:9}\par

\Quote{你要赶紧到我跟前来。因为德玛斯因爱了今世，就离弃了我，往得撒洛尼去了；克勒斯刻去了迦拉达，弟铎去了达尔玛提亚。只有路加同我在一起。你要带着马尔谷来，带他同来，因为他为服事我有益处。我派了提希苛去厄弗所。我在特洛阿卡颇留的那件外衣，你来的时候要带来，还有那些书籍，特别是那些羊皮卷。}\par

值得探究他为什么召唤弟茂德来他这里，因为弟茂德受托管理一个教会和整个民族。这不是出于傲慢。因为保禄本已准备到他那里去；因为我们见他说：\Quote{但若我迟延，你该知道在天主的家中应当怎样行事。}\BibleRef{弟前 3:15}但是他被一种强烈的必要性所阻止。他不再能自由行动。他在狱中，被尼禄拘禁，几乎濒临死亡。为了这事不在他见到门徒之前发生，因此他派人去叫他，渴望在死前见到他，并也许要托付许多事情给他。因此他说：\Quote{趁冬天以前赶快到我这里来。}\par

\Quote{因为德玛斯离弃了我，爱了今世。}他没有说\Quote{使我在离世前见到你}，这会让他忧伤，而是说\Quote{因为我独自一人}，他说，\Quote{没有一个人帮助或支持我。}\par

\Quote{因为德玛斯离弃了我，爱了今世，往得撒洛尼去了}；即他爱了自己的安逸和安全，宁愿在家享乐，也不愿与我同受苦难，分享我目前的危险。他责备他一人，并非出于责备，而是为了坚固我们，使我们不致因逃避劳苦和危险而软弱，因为这就是\Quote{爱了今世}。同时他也希望吸引他的门徒到他这里来。\par

\Quote{克勒斯刻去了迦拉达，弟铎去了达尔玛提亚。}\par

这些人他没有责备。因为弟铎是最令人钦佩的人之一，以致他把那岛的事务托付给他，我指的是那不小的岛，而是克里特的大岛。\par

\Quote{只有路加同我在一起。}因为他不可分离地依附于他。就是他写了福音和宗徒大事录；他献身于劳苦和学习，是个刚毅的人；保禄论到他写道：\Quote{他在福音上众教会都有称赞。}\BibleRef{格后 8:18}\par

\Quote{你要带着马尔谷来，带他同来，因为他为服事我有益处。}\par

他不是为了自己的安慰，而是为了福音的职务而需要他。因为即使被囚，他也没有停止宣讲。同样地，他派人去叫弟茂德，不是为了自己，而是为了福音的缘故，好使他的死不至于给信友们带来骚动，当他自己的许多门徒在场时，可以防止混乱，并安慰那些几乎无法承受他的死的人。因为在罗马的信徒很可能是有地位的人。\par

\Quote{我派了提希苛去厄弗所。我在特洛阿卡颇留的那件外衣，你来的时候要带来，还有那些书籍，特别是那些羊皮卷。}\par

这里译为\Quote{外衣}的词可能指一件衣服，或者有些人说是一个袋子，里面装着书籍。但他即将离世去见天主，还要书籍做什么呢？他很需要它们，好将它们托付给信友们，他们将会保存它们以代替他的教导。那么所有的信友都会受到很大打击，但尤其是那些在他死时在场并与他相处的人。但他要求那件外衣，是为了避免不得不接受别人的施舍。因为我们在别处看到他非常避免这一点；在别处，当他向从厄弗所来的人讲话时，他说：\Quote{你们知道，这两只手供应了我及同我一起的人的需用。}\BibleRef{宗 20:34}又说：\Quote{施予比领受更为有福。}\par

第14节。\Quote{铜匠亚历山大向我行了多恶；主要按他所行的报答他。}\BibleRef{弟后 4:14}\par

这里他再次提到他的受苦，并非仅仅为了责备和控告那人，而是为预备他的门徒面对冲突，使他能坚定地忍受。虽然引起这些磨难的是卑鄙可鄙、毫无荣誉的人，但他却说，这一切都应该勇敢地承受。因为凡是从大人物那里受屈的人，由于作恶者的尊贵，也得到不小的荣耀。但被卑鄙下贱的人伤害，则遭受更大的烦恼。\Quote{他向我行了多恶，}他说，即他在多方面迫害我。但这些事不会不受惩罚！因为主要按他的行为报答他。正如他上边说：\Quote{我受了何等磨难，但从这一切中，主都拯救了我。}\BibleRef{弟后 3:11}同样在这里，他通过双重考虑安慰他的门徒：即他本人受屈，而对方将因其恶行受报。这并不是说圣人们因迫害者的受罚而高兴，而是因为福音的事业需要这样，而且软弱者将从中得到安慰。\par

第15节。\Quote{你也要提防他，因为他极力反对我们的话。}\BibleRef{弟后 4:15}\par

这就是说，他对我们怀有敌意，并反对我们。他没有说：报复、惩罚、驱逐他，尽管因赐予他的恩宠他本可以这样做，但他没有这样做；他也没有武装弟茂德去对抗他，只命令他避开他，把报仇留给天主，而且为了安慰软弱者，他说了主要按行为报答他，这是一种预言而非咒骂。他之所以说这些，是为了预备他门徒的心灵，从下文也明显可见。但请看他是如何提到他其他的苦难的。\par

第16节。他说：\Quote{在我初次申诉时，没有一人站在我身边，众人都离弃了我；愿天主不将此罪归于他们。}\BibleRef{弟后 4:16}\par

你看见他如何饶恕他的朋友，尽管他们做了可悲的事？因为被外人轻视与被自己朋友轻视不是同一回事。你看见他极度的沮丧？不能说，我被外人攻击，却在朋友的关注与支持中得到安慰；因为连这些也出卖了我。他说：\Quote{众人都离弃了我}。这并非小过。因为在战场上抛弃暴露于危险中的人、退缩不敢面对敌人之手的人，理当被朋友打，因为完全出卖了他们的目标；在福音的事上更是如此。但他所说的\Quote{首次答辩}是什么？他曾站在尼禄面前，得以逃脱。但后来，因为他使尼禄的司酒官归化，就被斩首。这里又是对他门徒的鼓励，接下来。\par

17节。\BibleQuote{但主与我同在，并坚固了我} \BibleRef{弟后 4:17}\par

虽然被人离弃，天主却不允许他受任何伤害。他说：祂坚固了我，也就是说，祂赐我讲话的胆量。祂不容我沉沦。\par

\BibleQuote{使福音的宣讲借着我而完成}\par

即是说，得以完成。注意他的极大谦卑。他不是说祂坚固了我，因我配得祂的恩赐，而是说\Quote{宣讲}，即我所受托的，\Quote{得以被众人完全知晓}。就如有人穿着紫袍、戴着王冠，并因这情形而得保全。\par

\BibleQuote{并使万民都能听见}\par

这是什么？就是使福音的光辉和祂上智安排对我的眷顾，为众人所知。\par

\BibleQuote{我也从狮子口中被救了出来}\par

18节。\BibleQuote{主必要救我脱离各种凶恶的事} \BibleRef{弟后 4:18}\par

看他已多么接近死亡。他落入了狮子的嘴边。因尼禄的残暴及其统治的强横冒险，他称尼禄为狮子。他说：\Quote{主救了我，也必要救我}。但若他说\Quote{祂必要救我}，为何又说\Quote{我已准备被奠祭}？请注意措辞：他说\Quote{祂救我脱离了狮子的口}；而又说\Quote{祂必要救我}，不是从狮子的口，而是\Quote{脱离各种凶恶的事}。因为那时祂救我脱离了危险；但如今为福音已做了足够的事，祂还要再次救我脱离一切罪恶，也就是说，祂不容我带着定罪离去。因为他能够\BibleQuote{抵抗罪恶直到流血} \BibleRef{希 12:4} 而不屈服，这便是从另一只狮子，即魔鬼，得救，所以这保全比先前似乎被交出时的保全更大。\par

\BibleQuote{使我安全进入他天上的国；愿光荣归于他，于无穷世之世！阿们。}\par

这便是在那里发光时的救恩。但\Quote{祂要保全我进入祂的国}是什么意思？祂要救我脱离一切责罚，并在那里保全我。因为\BibleQuote{那在现世憎恨自己性命的，必要保存性命于永生} \BibleRef{若 12:25}。\par

\BibleQuote{愿光荣归于他}。看，这里是归给圣子的光荣颂。\par

19节。\BibleQuote{请问候普黎史拉和阿桂拉，以及敖乃息佛洛一家的人} \BibleRef{弟后 4:19}\par

因为他那时在罗马，关于这人他曾说：\BibleQuote{愿主在那一天赐他获得主的仁慈} \BibleRef{弟后 1:18}。通过提到他，他也使他的家人更热切于这样的善行。\par

\BibleQuote{请问候普黎史拉和阿桂拉}。这是那些他不断提及的人，也曾与他们同住，他们曾接待了阿颇罗。他先提起女人，我想是因为她更热心、更有信德，因为她曾接待了阿颇罗；但也可能是随意提及。这样的问候对他们而言并非微小的安慰；它表达了尊重与爱意，并分享了极大的恩宠。因为那位圣洁有福之人的单纯问候就足以充满领受者的恩宠。\par

20节。\BibleQuote{厄辣斯托仍留在格林多，但我把特洛斐摩留在米肋托，因为他病了} \BibleRef{弟后 4:20}\par

我们从\WorkTitle{宗徒大事录}得知，这位特洛斐摩和提希苛曾与他一同从犹太航行，并到处作他的同伴，也许因为他们比别人更热心。\par

\BibleQuote{我把特洛斐摩留在米肋托，因为他病了}。那你为何不治好他，反而留下他？宗徒并非凡事都能，或他们并非在所有场合都施行神恩，以免人们将他们看得过高。这同样可在他们之前那些有福的义人身上见到，例如梅瑟，他的声音微弱。为何这缺陷未被移除？不，他常被忧伤和沮丧所苦，且未能进入应许之地。\par

因为天主容许了许多事，为使人性的软弱得以显明。如果连这些有缺陷的麻木不仁的犹太人都能问：\Quote{那领我们出埃及地的梅瑟在哪里？}\BibleRef{出 32:1}那么，倘若梅瑟也领他们进入了应许之地，他们又岂会对他有别的态度？倘若他没有被允许被法郎的恐惧所压倒，他们岂不会把他当作神吗？我们看见吕斯特辣人对\Person{保禄}和\Person{巴尔纳伯}也是如此，以为他们是神，就撕裂衣服，跑到人群中喊着说：\Quote{诸位，你们为什么做这些事？我们也是和你们有同样性情的人。}\BibleRef{宗 14:14}再者，\Person{伯多禄}治好生来瘸腿的人后，当众人都惊奇于这奇迹时，他回答：\Quote{以色列人哪，为什么因这事而惊奇？或者为什么定睛看我们，好像凭我们自己的能力和圣德使这人行走呢？}\BibleRef{宗 3:12}听那有福的保禄也说：\Quote{有一根刺加在我身上，免得我过于高举自己。}\BibleRef{格后 12:7}但你说，这是谦卑的话。绝非如此。那刺并不是为了使他谦卑而加给他的，他这样说也不仅仅是出于谦卑。还可以为此提出其他的原因。因此，请留意天主如何为此解释，说：\Quote{我的恩宠为你足够了}，不是\Quote{免得你过于高举自己}，而是什么？\Quote{因为我的能力在软弱中才得以完全。}这样，两个目的同时达到了：所行的事清楚显明了，一切都归于天主。因此他在别处说：\Quote{我们将这宝贝放在瓦器里}\BibleRef{格后 4:7}，即在软弱、易受苦的身体里。为什么？\Quote{为要使这卓越的能力归于天主，而不是出于我们。}如果我们的身体不受软弱所困，一切都会归功于它们。我们也在别处看见他为\Person{厄帕洛狄托}的软弱而忧伤，关于他写道：\Quote{他病得几乎要死，但天主怜悯了他。}\BibleRef{斐 2:27}还有许多其他例子表明他对事件的不知情，这对他和他的门徒都是有益的。\par

\Quote{\Person{特洛斐摩}，我把他留在\Place{米肋托}了，他患病了。}\Place{米肋托}靠近\Place{厄弗所}。这事发生在他航行到犹太的时候，还是别的场合？因为他在罗马之后，又回到了\Place{西班牙}，但他是否再从那里来到这些地区，我们不知道。然而我们看见他被众人遗弃。\Quote{因为\Person{德玛斯}，}他说，\Quote{离弃了我。\Person{克勒斯刻}去了\Place{迦拉达}，\Person{弟铎}去了\Place{达尔玛提雅}。\Person{厄辣斯托}留在\Place{格林多}。\Person{特洛斐摩}，我把他留在\Place{米肋托}了，他患病了。}\par

第21节：\Quote{你要竭力争先在冬天以前来。\Person{欧步罗}问候你，\Person{普登}、\Person{利诺}、\Person{克劳狄雅}，以及众弟兄都问候你。}\BibleRef{弟后 4:21}\par

这位\Person{利诺}，有人说，是\Person{伯多禄}之后罗马教会的第二位主教。\Quote{以及\Person{克劳狄雅}。}你们看，妇女们多么热忱于信仰，多么炽烈！\Person{普黎史拉}就是这样，\Person{克劳狄雅}也是这样，她们已经钉在十字架上，已经准备好战斗！但为什么，既然有这么多信友，他却只提到这些妇女？显然是因为她们在心意上已经脱离了世俗事务，并且比其他人显赫。因为妇女作为女性，不会遇到任何障碍。这是天主恩宠的作为，使这一性别只在今世事务中受阻碍，或者即使在今世事务中也不受阻碍。因为妇女承担了整个家管中不小的份额，是家庭的守护者。没有她们，连政治事务也无法妥善进行。因为如果她们的家庭事务陷入混乱和无序，那些从事公共事务的人就会被拖在家里，政治事务就会处理不当。所以，即使在那些事务上，如同在属灵事务上一样，她也不逊色。因为她若能立志，就能忍受千死。因此，许多妇女殉道。她也能比男子更好地实践贞洁，因为没有如此强烈的火焰扰乱她；并能表现谦逊、庄重和\Quote{圣德，没有圣德，无人能见主}\BibleRef{希 12:14}；以及蔑视财富，只要她愿意，简言之，一切其他美德。\par

\Quote{你要竭力争先在冬天以前来。}看他如何催促他，却没有说任何使他伤心的话。他没有说\Quote{在我死以前}，免得使他忧伤；而是说\Quote{在冬天以前}，使你不被耽搁。\par

\Quote{\Person{欧步罗}，}他说，\Quote{问候你，\Person{普登}、\Person{利诺}、\Person{克劳狄雅}，以及众弟兄都问候你。}他没有提其余人的名字。你看见那些是最热忱的吗？\par

第22节：\Quote{愿主耶稣基督与你的心神同在。}\BibleRef{弟后 4:22}\par

没有比这更好的祈祷了。不要为我的离去而忧伤。主将与你同在。他说\Quote{与你}，而是\Quote{与你的心神同在}。因此有双重的助佑：圣神的恩宠，以及天主帮助它。否则，天主不会与我们同在，如果我们没有属神的恩宠。因为如果我们被恩宠抛弃，他又如何与我们同在呢？\par

\Quote{愿恩宠与你们同在。阿们。}\par

同样他也为自己祈祷，使他们常常中悦于他，使他们连同属神的恩赐一起拥有恩宠，因为哪里有了恩宠，就没有什么是令人忧伤的。因为正如那看到君王并蒙他宠爱的人，感觉不到任何不安；同样，即使我们的朋友离弃我们，即使我们遭遇灾难，如果那恩宠与我们同在并坚固我们，我们也不会感到痛苦。\par

道德方面。但我们怎样才能将恩宠吸引到我们身上呢？藉着做天主所喜悦的事，并在一切事上服从祂。在大户人家中，我们难道不是看见那些受宠爱的仆人，他们不关心自己的利益，反而以全部的热忱和敏捷来促进主人的利益，并且不是出于主人的强迫，而是出于自己的情感和良好的品性，将一切安排得井然有序？当他们常出现在主人眼前，当他们置身于家中，当他们不从事任何私事，也不为自己的事操心，而是将主人的事视为自己的事时，他们便受宠。因为那将自己的事视为主人之事的人，并没有真正把自己的事交给主人，而是使主人的利益成为自己的利益；他甚至在主人的事务中如同主人一样发号施令，与他平等地管理。他常常同样受到仆人们的畏惧，他所说的一切，主人也这样说，从此他被所有的敌人所惧怕。\par

如果那在世俗事务中优先考虑主人利益而非自己利益的人，并没有真正忽视自己的利益，反而更加促进了它；那么，在属灵的事上更是如此。轻视你自己的事，你将会领受天主的事。这是祂自己所愿意的。轻视每一个，并夺取天国。住在那里，而不是这里。在那里令人畏惧，而不是在这里。如果你在那里令人畏惧，你就不会畏惧人，而是畏惧魔鬼，甚至畏惧魔鬼本身。但如果你依赖属世的财富，你在他们眼中将是可鄙的，甚至也常常在人眼中如此。无论你的财富如何，你将是富足于奴役之物。但如果你轻视这些，你在王的家里将光芒四射。\par

宗徒们正是如此，他们轻视奴仆的房屋和世俗的财富！请看他们如何在主人的事务中发号施令。他们说：\Quote{让这个人从疾病中被解救，另一个人从魔鬼的附身中被解救；捆绑这个人，释放那个人。} 这些事是在地上由他们做的，却如同在天上成就。因为祂说：\Quote{凡你们在地上所捆绑的，在天上也要被捆绑。}\BibleRef{玛 18:18} 祂还赐给他们比祂自己的权柄更大的权柄。我并非说谎，这从祂的话中显明：\Quote{信我的人，要做比这些更大的事。}\BibleRef{若 14:12} 为什么如此？因为这份尊荣反射到主人身上。正如在我们自己的事务中，如果仆人拥有大权，主人就更加被钦佩，因为如果仆人如此有权能，那么命令他的主人更是如此。但如果有人忽略主人的服侍，只想着自己的妻子、儿子或仆人，并寻求富有，在那里积攒财宝，藉着偷窃和抢夺主人的财产，他立刻就败亡了，他的财富也与他一同灭亡。\par

因此，有了这些榜样，我恳求你们，不要在意我们的财产，好使我们自己受到关注；不，让我们轻视它们，好使我们获得它们。如果我们轻视它们，祂必会照顾它们；如果我们照顾它们，天主必会轻视它们。让我们劳碌于天主的事，而不是我们自己的事，或者更确切地说，真正劳碌于我们自己的事，因为祂的事就是我们的。我不是说天，也不是说地，也不是这世界的事物；这些都不配祂。它们同样属于信徒和不信者。那么，我说什么是祂的事呢？祂的荣耀和祂的国度。这些是属于祂的，并因祂的缘故也是我们的。如何呢？祂说：\Quote{如果我们与祂同死，也必与祂同活；如果我们受苦，也必与祂一同为王。}\BibleRef{弟后 2:11} 我们成了\Quote{同承继者}，并被称为祂的\Quote{弟兄}。为什么当我们下沉时，祂却拉我们上升？我们还要贫穷、乞丐多久？天堂已摆在我们面前，我们还留恋在地上吗？祂的国度已为我们敞开，我们还选择这世上的贫穷吗？不朽的生命已赐给我们，我们还为田地、木石而耗尽自己吗？要真正地富有。我希望你们如此。要贪婪、强取，我不责备你们。在这里，不贪婪是过错，不强取是该受责备的。那么，这是什么？\Quote{天国是以猛力夺取的，以猛力夺取的人就得了它。}\BibleRef{玛 11:12} 在那里，你要暴力！要强取！它不会因被夺取而减少。因为德行不会分裂，虔诚不会减少，天国也不会。当你夺取德行时，它反而增长；而属世的财物被夺取时则会减少。这从以下可见：在一座城中有一万人；如果所有人都夺取德行，它就增多，因为他们在一万件事上成为义人。如果没有人夺取它，它就减少，因为无处可寻。\par

由此可见，美善之物因多人拥有而倍增，但世俗财物反因攫取而减少。让我们不要满足于贫困，而要选择富足。天主正是在那些享受祂王国的人众多时才显得富足。\BibleQuote{因为祂是富足的}，经上说，\BibleQuote{对一切呼求祂的人}。\BibleRef{罗 10:12}那么，增加祂的财富罢；而你将通过占有它、贪恋它、强夺它来增加它。的确需要暴力。为何？因为有如此多的阻碍，如妻子儿女、忧虑和世俗事务；此外还有那些恶魔，以及统治它们的魔鬼。因此需要暴力，需要刚毅。那使用暴力的人必经受劳苦。如何？他忍受一切，与困境抗争。如何？他几乎尝试不可能的事。如果那些使用暴力的人尚且如此，而我们对可能的事都退缩，又怎能得胜？何时我们才能享受所追求的事物？\BibleQuote{勇猛的人}，经上说，\BibleQuote{夺取}天国\BibleQuote{以暴力}。需要暴力和强夺。因为天国并非简单地摆在我们面前，唾手可得。那用力夺取的人，总是清醒警醒，焦虑思虑，以便在合适时机下手。你们没看到在战争中，那些将要夺取战利品的人整夜警戒、全副武装吗？如果那些旨在夺取世俗财物的人整夜警醒、全副武装，那么我们这些希望夺取属灵事物的人，难道要在白日睡觉打鼾，始终赤身裸体、手无寸铁吗？因为陷于罪的人是赤手空拳的；而行义的人是武装的。我们不以施舍来坚固自己。我们不为自己准备燃烧的灯。我们不披戴属灵的甲胄。我们不学习通往那里的道路。我们不保持清醒警醒，因此无法夺取任何战利品。\par

如果想攻占一个王国，难道他不是将死亡以千种形式置于眼前吗？他不是全副武装，练习战术，以此目标做一切事，然后冲上前去攻击吗？但我们不这样做。我们想在睡觉时夺取战利品，因此空手而归。你们没看到掠夺者如何逃跑，如何快速移动吗？他们如何强行穿过一切？这里也需要敏捷。魔鬼在追逐你。他命令前面的人留住你。但如果你强壮，如果你警醒，你会踢开一个，推开另一个，像鸟一样逃脱一切。是的，如果你离开这里，如果你逃脱市场和喧闹——我指的是此生——并到达那些更高的领域，在未来世界。因为在那里，如同在旷野，没有喧闹，没有人打扰或阻止你的行程。\par

你已经夺取了？夺取之后还需要一点努力，以免所夺的被夺走。如果我们奔跑，如果我们不注视眼前任何事物，只考虑如何逃避那些阻碍我们的人，我们就能安全地保持所夺取的。你已经夺取了贞洁？不要停留；逃到魔鬼够不到的地方。如果他看到追不上你，就会停止追赶；正如我们，当看不到抢劫我们的人时，就会放弃追赶，不再追赶，也不叫别人捉贼，而是任凭他们逃走。所以，你要在开始时奋力奔跑，当你脱离魔鬼的魔掌后，他以后就不会攻击你，你将安全地享受那些不可言喻的福气。愿天主藉我们的主耶稣基督，使我们众人获得这些福气。愿光荣、权能、尊崇、崇拜归于祂，偕同圣父和圣神，从今直到永远，世世无尽。阿们。\par

\SourceRecord{Translated by Philip Schaff.；Nicene and Post-Nicene Fathers, First Series；Vol. 13.；Edited by Philip Schaff.；Buffalo, NY: Christian Literature Publishing Co.,；1889.；<http://www.newadvent.org/fathers/230710.htm>.}
